% ============================================================
%  PLAN DE TESIS  -  Documento maestro (estructura modular)
%
%  Escuela Profesional de Ingenieria Informatica y de Sistemas - UNSAAC
%  Titulo: Desarrollo de un sistema bidireccional de traduccion de voz
%          Quechua-Espanol mediante aprendizaje profundo para consultas
%          medicas en Cusco.
%
%  Estructura por capitulos (plantilla UNSAAC):
%    I.   Introduccion (sistema de 6 bloques)   -> 04AspectosGenerales
%    II.  Marco teorico y estado del arte        -> 05MarcoTeorico
%    III. Metodologia de la investigacion        -> 06Metodologia
%    IV.  Estudio de factibilidad                -> 07Factibilidad
%    V.   Aspectos administrativos               -> 08Administrativos
%    + Bibliografia y Anexos                      -> 09, 10
%
%  Compilar:  pdflatex main -> biber main -> pdflatex main -> pdflatex main
% ============================================================
\documentclass[12pt,a4paper]{article}
% ============================================================
%  00preambulo.tex  -  Preambulo consolidado del Plan de Tesis
%  (APA7 + espanol + graficos + matematica + matriz apaisada)
%  Compilar desde main.tex:  pdflatex -> biber -> pdflatex -> pdflatex
% ============================================================

% ---- Codificacion e idioma ----
\usepackage[utf8]{inputenc}
\usepackage[T1]{fontenc}
\usepackage{lmodern}
\usepackage[spanish,es-tabla]{babel}
\usepackage{csquotes}

% ---- Geometria (margenes tipo tesis) ----
\usepackage[a4paper,left=3cm,right=2.5cm,top=2.5cm,bottom=2.5cm]{geometry}

% ---- Interlineado 1.5 ----
\usepackage{setspace}
\onehalfspacing

% ---- Tipografia de titulos (se re-define mas abajo para APA7) ----
\usepackage{titlesec}

% ---- Tablas y figuras ----
\usepackage{graphicx}
\usepackage{booktabs}
\usepackage{array}
\usepackage{multirow}
\usepackage{tabularx}
\usepackage{float}
\usepackage[font=small,labelfont=bf]{caption}

% ---- Matematica (atencion, perdida, flujo de datos de la cascada) ----
\usepackage{amsmath}

% ---- Pagina apaisada para la matriz de consistencia ----
\usepackage{pdflscape}

% ---- Color de celdas / rampas ----
\usepackage{colortbl}

% ---- Graficos vectoriales (pgfplots / tikz) ----
\usepackage{pgfplots}
\pgfplotsset{compat=1.18}
\usepackage{tikz}
\usetikzlibrary{shapes.geometric,arrows.meta,positioning,calc}

% ---- Colores base ----
\usepackage{xcolor}
\definecolor{azulUNSAAC}{RGB}{0,71,133}
\definecolor{grisSuave}{RGB}{240,240,240}

% ---- Rutas de imagenes (logo, figuras) ----
\graphicspath{{imagenes/}{./}{../}}

% ---- Bibliografia APA7 (biblatex + biber) ----
\usepackage[
  backend=biber,
  style=apa,
  sortcites=true,
  sorting=nyt,
  language=spanish
]{biblatex}
\DeclareLanguageMapping{spanish}{spanish-apa}
% Fuentes: referencias generales (raiz del repo) y antecedentes de la revision.
\addbibresource{../referencias.bib}
\addbibresource{antecedentes.bib}

% ---- Listas e indices (tocloft ANTES de hyperref) ----
\usepackage{longtable}
\usepackage{algorithm}
\usepackage{algpseudocode}
\usepackage{tocloft}

% ---- Hipervinculos ----
\usepackage[hidelinks]{hyperref}
\hypersetup{
  colorlinks=true,
  linkcolor=black,
  citecolor=azulUNSAAC,
  urlcolor=azulUNSAAC
}

% ---- Icono ORCID ----
\usepackage{orcidlink}

% ---- Sangria de primer parrafo ----
\usepackage{indentfirst}

% ============================================================
%  Colores de la "ecuacion universal" del titulo (cafe/rojo/verde)
%  Esquema estandarizado (marco del docente): el titulo se descompone en tres
%  componentes, cada uno con un color fijo reutilizable en todo el proyecto.
% ============================================================
\definecolor{colProceso}{RGB}{122,66,20}   % cafe  -> Proceso / Accion
\definecolor{colMetodo}{RGB}{176,42,42}     % rojo  -> Tecnologia / Metodo
\definecolor{colContexto}{RGB}{31,110,54}   % verde -> Contexto / Aplicacion

% ============================================================
%  CORRECCIONES A OBSERVACIONES (revision del asesor)
% ============================================================
\definecolor{verdeOK}{RGB}{13,124,51}
% R1 consolidado: correcciones de la 1a ronda en NEGRO (texto normal).
\newcommand{\nuevo}[1]{#1}
\newenvironment{bloqueNuevo}{\par}{\par}
% R2 (2a ronda): aceptada; en NEGRO (texto normal).
\newcommand{\obsdos}[1]{#1}
\newenvironment{bloqueObsDos}{\par}{\par}
% R3: integrada al texto en NEGRO (versi{\'o}n para el jurado, sin resaltar correcciones).
\newcommand{\obstres}[1]{#1}
\newenvironment{bloqueObsTres}{\par}{\par}
% Mejoras propias de rigor (NO son una ronda del asesor): quedan en NEGRO (texto normal),
% integradas al documento como el resto. El macro se conserva por si luego se desea resaltar.
\newcommand{\obscuatro}[1]{#1}
\newenvironment{bloqueObsCuatro}{\par}{\par}

% ---- APA 7: numero y titulo ENCIMA de tablas y figuras; Nota al pie ----
\captionsetup{labelsep=newline,labelfont=bf,textfont=it,
              justification=raggedright,singlelinecheck=false}
\captionsetup[table]{position=above,skip=4pt}
\captionsetup[figure]{position=above,skip=6pt}
\newcommand{\notaAPA}[1]{\par\vspace{2pt}%
  \begin{minipage}{\linewidth}\footnotesize\textit{Nota.} #1\end{minipage}\par}

% ---- Estilo de tabla del proyecto (APA: cabecera negra sin relleno) ----
\newcommand{\hdr}[1]{\textbf{#1}}
\newcommand{\celnivel}[1]{\textbf{#1}}

% ---- Enlaces internos: siglas (quz/quy) y terminos del glosario ----
\DeclareRobustCommand{\sig}[1]{\hyperlink{sig:#1}{\texttt{#1}}}
\DeclareRobustCommand{\glo}[2]{\hyperlink{glo:#1}{#2}}
\newcommand{\destino}[1]{\hypertarget{#1}{}}

% ---- APA 7: encabezados ----
\titleformat{\section}[display]{\normalfont\Large\bfseries\filcenter}{Cap{\'i}tulo~\thesection}{0.3em}{}
\titleformat{name=\section,numberless}{\normalfont\Large\bfseries\filcenter}{}{0pt}{}
\titleformat{\subsection}{\normalfont\large\bfseries}{\thesubsection}{1em}{}
\titleformat{\subsubsection}{\normalfont\normalsize\bfseries\itshape}{\thesubsubsection}{1em}{}

% ---- Rampa estandar de los graficos: azul claro -> navy -> turquesa -> verde ----
\definecolor{g1}{HTML}{8AC4E6}
\definecolor{g2}{HTML}{66AEDA}
\definecolor{g3}{HTML}{4294CC}
\definecolor{g4}{HTML}{2E6DA0}
\definecolor{g5}{HTML}{1E4870}
\definecolor{g6}{HTML}{1F6476}
\definecolor{g7}{HTML}{229184}
\definecolor{g8}{HTML}{27A597}
\definecolor{g9}{HTML}{40B183}
\definecolor{g10}{HTML}{59BD6F}
\definecolor{g11}{HTML}{6AC362}
\definecolor{g12}{HTML}{7BC95A}

% ---- Numeracion por capitulos (I, II, ...) y subsecciones (2.1, 2.1.1) ----
\renewcommand{\thesection}{\Roman{section}}
\renewcommand{\thesubsection}{\arabic{section}.\arabic{subsection}}
\renewcommand{\thesubsubsection}{\arabic{section}.\arabic{subsection}.\arabic{subsubsection}}
\setcounter{secnumdepth}{3}
\setcounter{tocdepth}{3}

% ---- Nombres en espanol de los indices ----
\addto\captionsspanish{\renewcommand{\contentsname}{\'Indice general}}
\floatname{algorithm}{Algoritmo}
\renewcommand{\listalgorithmname}{\'Indice de algoritmos}

% ---- Indice de ecuaciones (lista propia via tocloft) ----
\newlistof{ecuaciones}{loe}{\'Indice de ecuaciones}
\newcommand{\regecuacion}[1]{%
  \addcontentsline{loe}{ecuaciones}{\protect\numberline{\theequation}#1}}
\let\cftecuacionesfont\cfttabfont
\let\cftecuacionespagefont\cfttabpagefont
\setlength{\cftecuacionesindent}{\cfttabindent}
\setlength{\cftecuacionesnumwidth}{\cfttabnumwidth}

% ---- Paginacion APA 7: numero en la esquina SUPERIOR DERECHA ----
\usepackage{fancyhdr}
\pagestyle{fancy}
\setlength{\headheight}{14.5pt}
\fancyhf{}
\fancyhead[R]{\thepage}
\renewcommand{\headrulewidth}{0pt}
\fancypagestyle{plain}{%
  \fancyhf{}%
  \fancyhead[R]{\thepage}%
  \renewcommand{\headrulewidth}{0pt}%
}


\begin{document}

% ---- Preliminares ----
% ================= CARÁTULA (plantilla UNSAAC) =================
% Sin ancla de p{\'a}gina en la car{\'a}tula para evitar el choque de destino
% "page.1" con la primera p{\'a}gina ar{\'a}biga del cuerpo.
\hypersetup{pageanchor=false}
\begin{titlepage}
\centering
\singlespacing

{\large UNIVERSIDAD NACIONAL DE SAN ANTONIO ABAD\\ DEL CUSCO\par}
\vspace{0.2cm}
{FACULTAD DE INGENIER{\'I}A EL{\'E}CTRICA, ELECTR{\'O}NICA, INFORM{\'A}TICA Y
M{\'E}CANICA\par}
{ESCUELA PROFESIONAL DE INGENIER{\'I}A INFORM{\'A}TICA Y DE SISTEMAS\par}
\vspace{0.6cm}

\IfFileExists{imagenes/logo_unsaac.png}{\includegraphics[width=4.6cm]{imagenes/logo_unsaac.png}}{%
  \IfFileExists{../logo_unsaac.png}{\includegraphics[width=4.6cm]{../logo_unsaac.png}}{}}
\vspace{0.6cm}

{\scshape Tesis\par}
\vspace{0.5cm}
\rule{\linewidth}{0.4pt}\par\vspace{0.3cm}
{\large\bfseries%
``\textcolor{colProceso}{DESARROLLO DE UN SISTEMA BIDIRECCIONAL DE TRADUCCI{\'O}N DE VOZ
QUECHUA--ESPA{\~N}OL}
\textcolor{colMetodo}{MEDIANTE APRENDIZAJE PROFUNDO}
\textcolor{colContexto}{PARA CONSULTAS M{\'E}DICAS EN CUSCO}''\par}
\vspace{0.3cm}\rule{\linewidth}{0.4pt}\par

\vspace{1.1cm}
\begin{flushleft}
\hspace*{0.28\textwidth}%
\begin{minipage}{0.66\textwidth}
\raggedright
{\small Para optar al t{\'i}tulo profesional de:}\\[3pt]
\hspace*{1.4em}INGENIERO INFORM{\'A}TICO Y DE SISTEMAS\\[13pt]
{\small Presentado por:}\\[3pt]
\hspace*{1.4em}Alain Armando Condori Contreras\\[13pt]
{\small C{\'o}digo de alumno:}\\[3pt]
\hspace*{1.4em}021501\\[13pt]
{\small Asesor:}\\[3pt]
\hspace*{1.4em}Mtro. Ing. Ra{\'u}l Huillca Huallparimachi
\end{minipage}
\end{flushleft}

\vfill
{\bfseries Per{\'u}, Noviembre de 2026\par}
\end{titlepage}
        % Caratula

% Preliminares en n{\'u}meros romanos (Resumen = p{\'a}gina I); la car{\'a}tula no se numera.
\pagenumbering{Roman}
\hypersetup{pageanchor=true}

% ================= RESUMEN =================
\section*{Resumen}
\addcontentsline{toc}{section}{Resumen}
\singlespacing\noindent
En la regi{\'o}n Cusco, m{\'a}s del 55\,\% de la poblaci{\'o}n tiene el quechua como lengua
materna y cerca de 70\,000 asegurados de EsSalud son quechuahablantes \parencite{essalud2023};
esta barrera idiom{\'a}tica afecta la calidad de la atenci{\'o}n en salud \parencite{begazo2026}.

El presente trabajo propone desarrollar y evaluar un sistema bidireccional de traducci{\'o}n de
voz Quechua--Espa{\~n}ol para consultas m{\'e}dicas, que encadena reconocimiento de voz (MMS),
traducci{\'o}n autom{\'a}tica (NLLB adaptado al quechua de Cusco, \sig{quz}) y s{\'i}ntesis de
voz (MMS-TTS) mediante el ajuste (\emph{fine-tuning}) de modelos preentrenados de c{\'o}digo
abierto.

La salida por voz en ambos sentidos busca incluir a la poblaci{\'o}n quechuahablante con
baja o nula alfabetizaci{\'o}n. La validaci{\'o}n se realizar{\'a} con consultas simuladas, sin datos
reales de pacientes, midiendo precisi{\'o}n (WER, BLEU, chrF), calidad de voz (MOS) y percepci{\'o}n
de utilidad. Se espera aportar el primer corpus cl{\'i}nico Quechua--Espa{\~n}ol y evidencia
emp{\'i}rica sobre la viabilidad de estas tecnolog{\'i}as en el dominio sanitario andino.

\vspace{0.3cm}
\textit{Palabras clave:} \nuevo{Traducci{\'o}n autom{\'a}tica, Quechua, Procesamiento del
lenguaje natural, Comunicaci{\'o}n m{\'e}dico-paciente, Inclusi{\'o}n digital.}
\onehalfspacing

\newpage
\section*{Abstract}
\addcontentsline{toc}{section}{Abstract}
\singlespacing\noindent
In the Cusco region, more than 55\% of the population has Quechua as their mother tongue and
nearly 70{,}000 EsSalud insured members are Quechua speakers \parencite{essalud2023}; this language
barrier affects the quality of health care \parencite{begazo2026}.

This work proposes to develop and
evaluate a bidirectional Quechua--Spanish speech translation \emph{system} for medical consultations,
chaining speech recognition (MMS), machine translation (NLLB adapted to Cusco Quechua, \sig{quz}) and
speech synthesis (MMS-TTS) through the \emph{fine-tuning} of open-source pretrained models.

The
two-way voice output aims to include Quechua-speaking people with low or no literacy. Validation will
be carried out with simulated consultations, without real patient data, measuring accuracy (WER,
BLEU, chrF), voice quality (MOS) and perceived usefulness. The study is expected to contribute the
first clinical Quechua--Spanish corpus and empirical evidence on the feasibility of these
technologies in the Andean health domain.

\vspace{0.3cm}
\noindent\obsdos{\textit{Keywords:} Machine translation, Quechua, Natural language processing,
Doctor--patient communication, Digital inclusion.}
\onehalfspacing
         % Resumen y Abstract
\newpage
\tableofcontents                                   % Índice general

\clearpage
\addcontentsline{toc}{section}{\'Indice de figuras}
\listoffigures                                     % Índice de figuras

\clearpage
\addcontentsline{toc}{section}{\'Indice de tablas}
\listoftables                                      % Índice de tablas

\clearpage
\addcontentsline{toc}{section}{\'Indice de algoritmos}
\listofalgorithms                                  % Índice de algoritmos

\clearpage                                          % Índice de símbolos
\section*{\'Indice de s{\'i}mbolos}
\addcontentsline{toc}{section}{\'Indice de s{\'i}mbolos}
\begin{longtable}{@{}p{2.6cm} p{9cm} p{2.4cm}@{}}
\toprule\hdr{S{\'i}mbolo} & \hdr{Descripci{\'o}n} & \hdr{\obstres{P{\'a}gina}} \\
\midrule \endhead
$\alpha$ & Nivel de significancia estad{\'i}stica & \obstres{\pageref{sym:alpha}} \\
$N$      & Tama{\~n}o de muestra / n.\ de palabras de referencia & \obstres{\pageref{eq:wer}} \\
$\beta$  & Factor de ponderaci{\'o}n recall/precisi{\'o}n en chrF & \obstres{\pageref{eq:chrf}} \\
$S,D,I$  & Sustituciones, eliminaciones e inserciones (c{\'a}lculo del WER) & \obstres{\pageref{eq:wer}} \\
$\theta$ & Par{\'a}metros (pesos) del modelo ajustado por \emph{fine-tuning} & \obstres{\pageref{eq:perdida}} \\
$\mathcal{L}$ & Funci{\'o}n de p{\'e}rdida (entrop{\'i}a cruzada) del entrenamiento & \obstres{\pageref{eq:perdida}} \\
\bottomrule
\end{longtable}

\clearpage                                          % Índice de abreviaturas
\section*{\'Indice de abreviaturas}
\addcontentsline{toc}{section}{\'Indice de abreviaturas}
\begin{longtable}{@{}p{3cm} p{12cm}@{}}
\toprule\hdr{Sigla} & \hdr{Significado} \\
\midrule \endhead
IA      & Inteligencia artificial \\
PLN     & Procesamiento del lenguaje natural \\
ASR     & Reconocimiento autom{\'a}tico del habla (\emph{Automatic Speech Recognition}) \\
STT     & Voz a texto (\emph{Speech-to-Text}) \\
TTS     & S{\'i}ntesis de voz (\emph{Text-to-Speech}) \\
MT      & Traducci{\'o}n autom{\'a}tica (\emph{Machine Translation}) \\
MMS     & \emph{Massively Multilingual Speech} (Meta) \\
NLLB    & \emph{No Language Left Behind} (Meta) \\
WER     & Tasa de error de palabra (\emph{Word Error Rate}) \\
BLEU    & \emph{Bilingual Evaluation Understudy} \\
chrF    & \emph{Character n-gram F-score} \\
MOS     & \emph{Mean Opinion Score} \\
quz     & Quechua Collao / Cusco (c{\'o}digo ISO 639-3)\destino{sig:quz} \\
quy     & Quechua Ayacucho (c{\'o}digo ISO 639-3)\destino{sig:quy} \\
Transformer & \nuevo{Arquitectura de red neuronal basada en autoatenci{\'o}n (base de MMS/NLLB)} \\
INEI    & Instituto Nacional de Estad{\'i}stica e Inform{\'a}tica \\
EsSalud & Seguro Social de Salud del Per{\'u} \\
UNSAAC  & Universidad Nacional de San Antonio Abad del Cusco \\
\bottomrule
\end{longtable}

\clearpage                                          % Índice de ecuaciones
\section*{\'Indice de ecuaciones}
\addcontentsline{toc}{section}{\'Indice de ecuaciones}
\makeatletter\@starttoc{loe}\makeatother

\clearpage                                          % Glosario de términos
\section*{Glosario de t{\'e}rminos}
\addcontentsline{toc}{section}{Glosario de t{\'e}rminos}
\begin{longtable}{@{}p{3.6cm} p{11.4cm}@{}}
\toprule\hdr{T{\'e}rmino} & \hdr{Definici{\'o}n} \\
\midrule \endhead
Corpus                  & Conjunto estructurado de textos o grabaciones usado para entrenar o evaluar modelos.\destino{glo:corpus} \\
Fine-tuning             & Ajuste de un modelo preentrenado con datos de un dominio espec{\'i}fico.\destino{glo:finetuning} \\
Lengua de bajo recurso  & Idioma con escasez de datos digitales para entrenar modelos de PLN.\destino{glo:bajorecurso} \\
Modelo preentrenado     & Modelo entrenado sobre grandes vol{\'u}menes de datos, reutilizable mediante ajuste.\destino{glo:preentrenado} \\
Lengua polisint{\'e}tica & Lengua que concentra mucha informaci{\'o}n gramatical en una sola palabra, como el quechua.\destino{glo:polisintetica} \\
Pipeline                & Secuencia encadenada de etapas (voz $\to$ texto $\to$ traducci{\'o}n $\to$ voz).\destino{glo:pipeline} \\
\bottomrule
\end{longtable}

\clearpage
\setcounter{table}{0}\setcounter{figure}{0}
% Inicio del cuerpo: paginaci{\'o}n normal en ar{\'a}bigos desde 1 (APA).
\pagenumbering{arabic}
         % Indices, glosario, abreviaturas

% ---- Cuerpo ----
% ##################################################################
% CAPÍTULO I: INTRODUCCIÓN
% ##################################################################
\section{Introducci{\'o}n}
En la regi{\'o}n Cusco, donde m{\'a}s de la mitad de la poblaci{\'o}n tiene el quechua como lengua
materna, la atenci{\'o}n en salud opera de manera predominante en espa{\~n}ol. Esta asimetr{\'i}a
levanta una barrera comunicativa entre un personal de salud mayoritariamente castellanohablante y una
poblaci{\'o}n usuaria quechuahablante, barrera que ---agravada por la baja alfabetizaci{\'o}n en las
zonas rurales andinas--- compromete aspectos sensibles del acto m{\'e}dico, desde la recogida de los
s{\'i}ntomas hasta la correcta comprensi{\'o}n de las indicaciones.

En paralelo, el procesamiento del lenguaje natural ha alcanzado un punto en el que es posible
\emph{adaptar} grandes modelos multiling{\"u}es preentrenados a lenguas de bajo recurso como el
quechua, sin necesidad de entrenarlos desde cero. Sobre esa base, el presente plan propone desarrollar
y evaluar un sistema \textbf{bidireccional de traducci{\'o}n de voz Quechua--Espa{\~n}ol} para
consultas m{\'e}dicas, que encadena reconocimiento de voz, traducci{\'o}n autom{\'a}tica y s{\'i}ntesis
de voz, con salida hablada en ambos sentidos del di{\'a}logo para incluir tambi{\'e}n a los usuarios
con baja o nula alfabetizaci{\'o}n.

Este primer cap{\'i}tulo sit{\'u}a el problema en sus niveles internacional, nacional y regional,
delimita el vac{\'i}o de investigaci{\'o}n y el problema espec{\'i}fico, y formula la propuesta con sus
objetivos, su justificaci{\'o}n y sus alcances. Las hip{\'o}tesis, por su car{\'a}cter contrastable, se
enuncian en el Cap{\'i}tulo~III junto con el dise{\~n}o metodol{\'o}gico que permite someterlas a
prueba.

\subsection{Contexto y problema general}
\obsdos{El contexto de esta investigaci{\'o}n se aborda en tres niveles ---internacional, nacional y
regional--- que, en conjunto, muestran c{\'o}mo una barrera ling{\"u}{\'i}stica hist{\'o}rica se
combina hoy con una oportunidad tecnol{\'o}gica inexistente hace pocos a{\~n}os.}

\subsubsection{\obsdos{Contexto internacional}}
\begin{bloqueObsDos}
\noindent De las miles de lenguas que se hablan en el mundo, solo una peque{\~n}a fracci{\'o}n cuenta
con recursos digitales suficientes para sostener tecnolog{\'i}as del lenguaje; la gran mayor{\'i}a
son \emph{lenguas de bajo recurso}, con escasa o nula presencia en los sistemas de voz y
traducci{\'o}n autom{\'a}tica \parencite{joshi2020diversity}. Esta asimetr{\'i}a digital reproduce, en
el plano tecnol{\'o}gico, las desigualdades que ya enfrentan sus hablantes en el acceso a servicios
b{\'a}sicos.

\noindent En el {\'a}mbito
sanitario, la comunicaci{\'o}n entre un profesional de salud y un paciente que no comparten la misma
lengua es un problema reconocido a escala global: la barrera idiom{\'a}tica compromete la
anamnesis, el diagn{\'o}stico, el consentimiento informado y la adherencia al tratamiento
\parencite{alshamsi2020barreras}, y por ello el derecho de la persona a ser atendida en su lengua
materna se considera un componente de la calidad y la seguridad del acto m{\'e}dico.

\noindent En paralelo, el procesamiento del lenguaje natural ha experimentado un cambio de
paradigma que abre una ventana de oportunidad para estas lenguas. La estrategia dominante ya no es
entrenar modelos desde cero ---inviable con pocos datos---, sino adaptar (\emph{fine-tuning})
grandes modelos multiling{\"u}es preentrenados que incorporan cientos de idiomas: en la voz, MMS
\parencite{pratap2024mms} y Whisper \parencite{radford2023whisper}; en la traducci{\'o}n, NLLB
\parencite{costajussa2022nllb}.

\noindent Competencias internacionales como AmericasNLP
\parencite{mager2021americasnlp} y las tareas de traducci{\'o}n de voz de IWSLT
\parencite{Ortega2025260} han impulsado de forma sostenida el trabajo sobre lenguas ind{\'i}genas
americanas. En este panorama, el quechua ---la lengua originaria con mayor n{\'u}mero de hablantes
de Am{\'e}rica, extendida por varios pa{\'i}ses andinos--- concentra buena parte de ese esfuerzo,
aunque todav{\'i}a de manera fragmentaria y por componentes aislados.

\noindent En el terreno sanitario, la pr{\'a}ctica internacional recomendada frente a la barrera
idiom{\'a}tica es la mediaci{\'o}n ling{\"u}{\'i}stica ---int{\'e}rpretes presenciales o
telef{\'o}nicos--- \parencite{alshamsi2020barreras}; sin embargo, su disponibilidad es limitada y
pr{\'a}cticamente nula en entornos rurales y de atenci{\'o}n primaria, donde ocurre el primer
contacto con el sistema de salud.

\noindent Por eso
las tecnolog{\'i}as del habla ---reconocimiento, traducci{\'o}n y s{\'i}ntesis de voz--- se perfilan
a escala global como un apoyo emergente a la comunicaci{\'o}n cl{\'i}nica, con la ventaja decisiva de
operar por \emph{voz}: a diferencia de las soluciones basadas en texto, no exigen alfabetizaci{\'o}n
y resultan utilizables por poblaciones con baja escolaridad.

\noindent Esta orientaci{\'o}n se inscribe en un
esfuerzo m{\'a}s amplio de \emph{inclusi{\'o}n digital} y de preservaci{\'o}n de las lenguas
originarias, hoy reconocido como un objetivo de equidad tecnol{\'o}gica: dotar de presencia digital a
una lengua es tambi{\'e}n una forma de sostener los derechos de sus hablantes.
\end{bloqueObsDos}

\subsubsection{\obsdos{Contexto nacional}}
El Per{\'u} es un pa{\'i}s plurilingüe en el que el quechua es la lengua originaria con mayor
n{\'u}mero de hablantes. Los Censos Nacionales 2017 registran a millones de peruanos que declaran
el quechua como lengua materna, con una marcada concentraci{\'o}n en la sierra sur
\parencite{inei2017}.
\begin{bloqueObsDos}
\noindent El Estado reconoce el car{\'a}cter multiling{\"u}e del pa{\'i}s y el derecho de los
ciudadanos a relacionarse con las instituciones p{\'u}blicas en su lengua materna; sin embargo, en
la pr{\'a}ctica los servicios esenciales ---la salud entre ellos--- operan de manera predominante en
espa{\~n}ol, lo que condiciona su acceso efectivo para la poblaci{\'o}n quechuahablante.

\noindent A esta
brecha ling{\"u}{\'i}stica se suma una brecha educativa: cerca del 22.8\,\% de las mujeres rurales no
sabe leer ni escribir \parencite{inei2023}, con programas de alfabetizaci{\'o}n a{\'u}n limitados en
las zonas andinas \parencite{dispurse2023}.

\noindent Esta doble barrera ---de idioma y de
alfabetizaci{\'o}n--- implica que una soluci{\'o}n tecnol{\'o}gica pertinente para el pa{\'i}s no
puede basarse solo en texto, sino que debe operar principalmente por \emph{voz}. En el plano
institucional ya existen se{\~n}ales de respuesta: EsSalud reactiv{\'o} en 2023 una central
telef{\'o}nica de atenci{\'o}n en quechua \parencite{essalud2023}, y la comunidad acad{\'e}mica
peruana ha producido recursos y modelos iniciales para el quechua en reconocimiento, traducci{\'o}n
y s{\'i}ntesis de voz, aunque ninguno orientado a la consulta m{\'e}dica.

\noindent El sistema de salud peruano incorpora, al menos en el plano normativo, un enfoque
intercultural que reconoce la pertinencia ling{\"u}{\'i}stica de la atenci{\'o}n
\parencite{minsa2016intercultural}; no obstante, su
traducci{\'o}n en herramientas operativas dentro del consultorio sigue siendo escasa. La atenci{\'o}n
primaria se presta en una red amplia y dispersa de establecimientos ---postas y centros de salud---
situados en zonas andinas donde la poblaci{\'o}n usuaria es mayoritariamente quechuahablante, lo que
convierte la distancia idiom{\'a}tica en un problema cotidiano y estructural, no excepcional.

\noindent En ese
contexto, la investigaci{\'o}n nacional en procesamiento del lenguaje para el quechua ha crecido de
forma sostenida en la {\'u}ltima d{\'e}cada, pero se ha concentrado en resolver cada componente por
separado ---reconocimiento, traducci{\'o}n o s{\'i}ntesis--- y sobre dominios generales, sin integrar
un sistema completo ni abordar el vocabulario propio de la atenci{\'o}n en salud.
\end{bloqueObsDos}

\subsubsection{\obsdos{Contexto regional (Cusco)}}
En la regi{\'o}n Cusco la presencia del quechua es particularmente intensa: m{\'a}s del
55\,\% de la poblaci{\'o}n tiene el quechua como lengua materna \parencite{inei2017}. Esta
realidad ling{\"u}{\'i}stica condiciona el acceso efectivo a los servicios p{\'u}blicos
esenciales ---la atenci{\'o}n en salud entre ellos--- cuando operan de manera predominante en
espa{\~n}ol.

\begin{bloqueObsDos}
\noindent La red de salud de Cusco atiende a una poblaci{\'o}n extensa y geogr{\'a}ficamente
dispersa, con numerosas comunidades rurales y altoandinas donde el quechua es la lengua de uso
cotidiano y el espa{\~n}ol tiene presencia limitada.

\noindent A esta realidad se suma una particularidad
ling{\"u}{\'i}stica decisiva para la propuesta: el quechua de Cusco corresponde a la variante Collao
(\sig{quz}), distinta de la de Ayacucho (\sig{quy}) tanto en l{\'e}xico como en rasgos
fonol{\'o}gicos; esta diferencia dialectal ---no menor para un sistema autom{\'a}tico--- es
justamente la que los modelos de traducci{\'o}n disponibles a{\'u}n no cubren para Cusco.

\noindent Por
a{\~n}adidura, la consulta m{\'e}dica no es solo un intercambio de informaci{\'o}n, sino un encuentro
cultural en el que la confianza y la comprensi{\'o}n mutua condicionan la adherencia al tratamiento;
cuando esa comunicaci{\'o}n falla, el riesgo cl{\'i}nico recae sobre la poblaci{\'o}n m{\'a}s
vulnerable.
\end{bloqueObsDos}

\subsubsection{\obsdos{Planteamiento del problema}}
\obsdos{Los tres niveles de contexto ---internacional, nacional y regional--- convergen en un
problema concreto y localizado: la comunicaci{\'o}n cl{\'i}nica en la consulta cara a cara de los
establecimientos de salud de Cusco, donde el idioma se vuelve un determinante directo de la
calidad y la seguridad de la atenci{\'o}n.}

En el {\'a}mbito sanitario, la distancia idiom{\'a}tica entre un personal de salud mayoritariamente
castellanohablante y una poblaci{\'o}n usuaria quechuahablante origina una barrera comunicativa con
consecuencias cl{\'i}nicas directas.

EsSalud reconoce que cerca de 70\,000 de sus asegurados en
Cusco son quechuahablantes y reactiv{\'o} en 2023 la central telef{\'o}nica en quechua
\parencite{essalud2023}; sin embargo, \obsdos{esa medida ---de alcance telef{\'o}nico y
unidireccional--- no alcanza} la interacci{\'o}n cara a cara del
consultorio, donde se recogen s{\'i}ntomas, se formula el diagn{\'o}stico y se entregan las
indicaciones.

La evidencia confirma la gravedad del problema: \textcite{begazo2026} documenta
c{\'o}mo la desconexi{\'o}n idiom{\'a}tica vulnera la relaci{\'o}n terap{\'e}utica en hospitales de
Cusco, y \textcite{cuaila2022} halla una relaci{\'o}n entre el conocimiento del quechua del personal
y la calidad de atenci{\'o}n percibida. A ello se suma una barrera de alfabetizaci{\'o}n: cerca del
22.8\,\% de las mujeres rurales no sabe leer ni escribir \parencite{inei2023}, lo que obliga a una
soluci{\'o}n basada en \emph{voz} y no solo en texto (Figura~\ref{fig:barreras}).

\obsdos{\noindent Las consecuencias de esta barrera no son menores: una anamnesis incompleta, un
diagn{\'o}stico tard{\'i}o o una indicaci{\'o}n mal comprendida comprometen la seguridad del acto
m{\'e}dico y la adherencia al tratamiento, y recaen sobre la poblaci{\'o}n m{\'a}s vulnerable. Pese a
ello, las respuestas institucionales vigentes son parciales ---telef{\'o}nicas, unidireccionales o
dependientes de la disponibilidad de personal biling{\"u}e--- y no existe una herramienta que asista
la traducci{\'o}n por voz, en ambos sentidos, durante la propia consulta.}

\begin{bloqueObsDos}
\noindent En s{\'i}ntesis, el problema general que motiva este plan es la \textbf{barrera
comunicativa entre el personal de salud castellanohablante y los pacientes quechuahablantes de
Cusco} durante la consulta cara a cara, agravada por la baja alfabetizaci{\'o}n y no resuelta por
las medidas vigentes. Existiendo hoy modelos multiling{\"u}es preentrenados que cubren parcialmente
el quechua, la pregunta es si pueden adaptarse e integrarse en un sistema de traducci{\'o}n por voz,
bidireccional, para asistir esa interacci{\'o}n cl{\'i}nica.
\end{bloqueObsDos}

\begin{figure}[H]
\centering
\caption{Magnitud de la problem{\'a}tica en Cusco.}
\label{fig:barreras}
\begin{tikzpicture}
\begin{axis}[
    ybar, width=12cm, height=6cm, bar width=28pt,
    ymin=0, ymax=100, ylabel={Porcentaje (\%)},
    symbolic x coords={Lengua materna,Asegurados EsSalud,Analfab. mujer rural},
    xtick={Lengua materna,Asegurados EsSalud,Analfab. mujer rural},
    nodes near coords, nodes near coords style={font=\small},
    x tick label style={align=center,font=\small,text width=2.5cm},
    enlarge x limits=0.25, ymajorgrids=true, grid style=dashed,
]
\addplot[fill=g4, draw=g5, bar shift=0pt] coordinates {(Lengua materna,55)};
\addplot[fill=g8, draw=g7, bar shift=0pt] coordinates {(Asegurados EsSalud,15)};
\addplot[fill=g12, draw=g11, bar shift=0pt] coordinates {(Analfab. mujer rural,22.8)};
\end{axis}
\end{tikzpicture}
\notaAPA{Elaboraci{\'o}n propia con datos de \textcite{inei2017}, \textcite{essalud2023} e \textcite{inei2023}.}
\end{figure}

\subsection{Tecnolog{\'i}a o m{\'e}todo}
El estado del arte en procesamiento del lenguaje natural ofrece hoy una oportunidad que no
exist{\'i}a hace pocos a{\~n}os. La estrategia dominante para lenguas de bajo recurso ya no consiste
en entrenar modelos desde cero, sino en adaptar (\emph{fine-tuning}) grandes modelos multiling{\"u}es
preentrenados.

En la voz, el proyecto MMS \parencite{pratap2024mms} cubre el quechua de Cusco
(\sig{quz}) en reconocimiento y s{\'i}ntesis; Whisper \parencite{radford2023whisper} declara una
cobertura marginal del quechua y se reserva como l{\'i}nea base. En traducci{\'o}n, NLLB
\parencite{costajussa2022nllb} incluye el quechua de Ayacucho (\sig{quy}) y no el de Cusco, por lo
que el par debe adaptarse a \sig{quz} ---esa brecha dialectal es el n{\'u}cleo del aporte
t{\'e}cnico---.

Existen adem{\'a}s antecedentes de traducci{\'o}n neural para esta lengua
polisint{\'e}tica \parencite{ortega2020quechua,mager2021americasnlp} y recursos como el corpus
\emph{Huqariq} \parencite{zevallos2022huqariq}.

\subsection{Vac{\'i}o de investigaci{\'o}n}
Persiste una desconexi{\'o}n entre ambos mundos: el problema social est{\'a} bien documentado y la
tecnolog{\'i}a base existe, pero nadie ha unido los dos extremos. No se dispone de un corpus
cl{\'i}nico Quechua--Espa{\~n}ol ni de una evaluaci{\'o}n emp{\'i}rica del desempe{\~n}o de estos
modelos sobre el vocabulario propio de la consulta m{\'e}dica (triaje, s{\'i}ntomas e indicaciones),
y NLLB a{\'u}n no cubre la variante de Cusco (\sig{quz}). Cerrar esa intersecci{\'o}n es el aporte
de este proyecto.

\subsection{Formulaci{\'o}n del problema}
La formulaci{\'o}n del problema de esta investigaci{\'o}n se concentra en la comunicaci{\'o}n
m{\'e}dico-paciente mediada por un sistema bidireccional de traducci{\'o}n de voz Quechua--Espa{\~n}ol
basado en modelos preentrenados, y en la medida en que su ajuste al dominio cl{\'i}nico mejora la
precisi{\'o}n y la comprensi{\'o}n frente a la ausencia de apoyo tecnol{\'o}gico.

\subsubsection{Problema general}
{\textquestiondown}\obsdos{C{\'o}mo \emph{desarrollar un sistema bidireccional de traducci{\'o}n de voz
Quechua--Espa{\~n}ol} \emph{mediante arquitecturas de transformadores multimodales} (ajuste fino de los
modelos preentrenados MMS y NLLB) \emph{para la comunicaci{\'o}n m{\'e}dico-paciente en las consultas m{\'e}dicas en Cusco}?}

\subsubsection{Problemas espec{\'i}ficos}
\begin{enumerate}
  \item \obsdos{\textbf{(PE1 --- Diagn{\'o}stico).} {\textquestiondown}Cu{\'a}les son las barreras de comunicaci{\'o}n entre el personal de salud y los pacientes quechuahablantes, y qu{\'e} debe contener el corpus cl{\'i}nico Quechua--Espa{\~n}ol (triaje, s{\'i}ntomas e indicaciones) para poder ajustar los modelos? \textbf{Aporte:} Permite construir el corpus cl{\'i}nico etiquetado; sin esos datos no es posible ajustar los modelos.}
  \item \obsdos{\textbf{(PE2 --- Construcci{\'o}n).} {\textquestiondown}C{\'o}mo construir el sistema de traducci{\'o}n por voz, ajustando (\emph{fine-tuning}) los modelos preentrenados MMS y NLLB ---del quechua de Ayacucho (\sig{quy}) al de Cusco (\sig{quz})--- e integr{\'a}ndolos en un \emph{pipeline} ASR+MT+TTS? \textbf{Aporte:} Es el n{\'u}cleo de la soluci{\'o}n: el sistema con el modelo de traducci{\'o}n adaptado al dominio cl{\'i}nico y a la variante de Cusco.}
  \item \obsdos{\textbf{(PE3 --- Validaci{\'o}n).} {\textquestiondown}En qu{\'e} medida el sistema mejora la precisi{\'o}n de la traducci{\'o}n (WER, BLEU, chrF), la calidad de la voz (MOS) y la comprensi{\'o}n entre el m{\'e}dico y el paciente, frente a no contar con apoyo tecnol{\'o}gico? \textbf{Aporte:} Demuestra si el sistema resuelve el problema general.}
\end{enumerate}

\subsection{Propuesta de soluci{\'o}n}
Se propone desarrollar y evaluar un sistema bidireccional de traducci{\'o}n de voz
Quechua--Espa{\~n}ol para consultas m{\'e}dicas, que encadena reconocimiento de voz (MMS),
traducci{\'o}n autom{\'a}tica (NLLB adaptado de \sig{quy} a \sig{quz}) y s{\'i}ntesis de voz
(MMS-TTS) mediante el ajuste de modelos preentrenados de c{\'o}digo abierto, con salida por voz en
ambos sentidos del di{\'a}logo.

\subsection{Objetivos}

\subsubsection{Objetivo general}
\nuevo{Desarrollar un sistema bidireccional de traducci{\'o}n de voz Quechua--Espa{\~n}ol para la
comunicaci{\'o}n m{\'e}dico-paciente, mediante arquitecturas de transformadores multimodales (ajuste fino
de los modelos preentrenados MMS y NLLB), en las consultas m{\'e}dicas en Cusco.}

\subsubsection{Objetivos espec{\'i}ficos}
\begin{enumerate}
  \item \obsdos{\textbf{(OE1 --- Diagn{\'o}stico).} Caracterizar las barreras de comunicaci{\'o}n entre el personal de salud y los pacientes quechuahablantes y determinar qu{\'e} debe contener el corpus cl{\'i}nico Quechua--Espa{\~n}ol (triaje, s{\'i}ntomas e indicaciones) para poder ajustar los modelos.}
  \item \obsdos{\textbf{(OE2 --- Construcci{\'o}n).} Construir el sistema de traducci{\'o}n por voz, ajustando (\emph{fine-tuning}) los modelos preentrenados MMS y NLLB ---del quechua de Ayacucho (\sig{quy}) al de Cusco (\sig{quz})--- e integr{\'a}ndolos en un \emph{pipeline} ASR+MT+TTS con salida por voz y texto.}
  \item \obsdos{\textbf{(OE3 --- Validaci{\'o}n).} Evaluar en qu{\'e} medida el sistema mejora la precisi{\'o}n de la traducci{\'o}n (WER, BLEU, chrF), la calidad de la voz (MOS) y la comprensi{\'o}n entre el m{\'e}dico y el paciente, frente a no contar con apoyo tecnol{\'o}gico.}
\end{enumerate}

\subsection{Justificaci{\'o}n}
\noindent La investigaci{\'o}n se justifica en cuatro dimensiones:
\begin{itemize}\setlength{\itemsep}{1pt}
  \item \textbf{Te{\'o}rica:} Aporta evidencia emp{\'i}rica del desempe{\~n}o de modelos
  preentrenados en una lengua de bajo recurso y morfolog{\'i}a compleja.
  \item \textbf{Pr{\'a}ctica:} Entrega un prototipo piloteable en un centro de salud.
  \item \textbf{Social:} Ataca una inequidad documentada en el acceso a la salud
  \parencite{begazo2026}, en l{\'i}nea con el derecho de los pacientes a ser atendidos en su
  lengua materna \parencite{essalud2023}.
  \item \textbf{Metodol{\'o}gica:} Establece un proceso replicable (corpus + ajuste + evaluaci{\'o}n)
  aplicable a otras variantes del quechua o a otros dominios.
\end{itemize}
\noindent Entre las \textbf{contribuciones originales esperadas} destacan:
\begin{itemize}\setlength{\itemsep}{1pt}
  \item El \textbf{primer corpus de dominio m{\'e}dico} Quechua--Espa{\~n}ol para tecnolog{\'i}as del habla.
  \item La \textbf{adaptaci{\'o}n de NLLB de \sig{quy} a \sig{quz}} en el dominio cl{\'i}nico
  (aporte t{\'e}cnico central).
  \item Un \textbf{sistema integrado ASR+MT+TTS} con salida por voz para usuarios con baja alfabetizaci{\'o}n.
\end{itemize}
\noindent En conjunto, el proyecto contribuye a la equidad en el acceso a la salud de la
poblaci{\'o}n quechuahablante de Cusco, a la seguridad del acto m{\'e}dico y a la revitalizaci{\'o}n
y el prestigio digital del quechua.

\subsubsection{Alcances y limitaciones}
\noindent\textbf{Alcances.}
\begin{itemize}
  \item Par ling{\"u}{\'i}stico Quechua de Cusco (\sig{quz}) -- Espa{\~n}ol, en ambos sentidos.
  \item Vocabulario acotado de consulta general (triaje, s{\'i}ntomas, indicaciones).
  \item Prototipo integrado ASR+MT+TTS con salida por voz y texto.
  \item Validaci{\'o}n con consultas simuladas en un centro de salud de referencia en Cusco.
\end{itemize}
\noindent\textbf{Limitaciones.}
\begin{itemize}
  \item \textbf{Temporal:} Ejecuci{\'o}n entre octubre 2026 y septiembre 2027 (12 meses).
  \item \textbf{{\'E}tica:} Se emplean di{\'a}logos simulados/actuados, sin datos reales de pacientes.
  \item Desempe{\~n}o condicionado por la escasez de datos de dominio cl{\'i}nico en quechua.
  \item Resultados circunscritos a la variante de Cusco; no generalizables sin ajuste a otras variantes.
\end{itemize}

\subsection{Organizaci{\'o}n de la tesis}
El documento se organiza en cinco cap{\'i}tulos:
\begin{itemize}\setlength{\itemsep}{1pt}
  \item \textbf{Cap{\'i}tulo~I --- Introducci{\'o}n:} contexto, tecnolog{\'i}a, vac{\'i}o de
  investigaci{\'o}n, problema, propuesta, objetivos y justificaci{\'o}n.
  \item \textbf{Cap{\'i}tulo~II --- Marco te{\'o}rico:} fundamentos te{\'o}ricos y estado del arte
  derivado de la revisi{\'o}n sistem{\'a}tica.
  \item \textbf{Cap{\'i}tulo~III --- Metodolog{\'i}a:} enfoque, dise{\~n}o, \obsdos{hip{\'o}tesis},
  variables, unidades de an{\'a}lisis, t{\'e}cnicas, procedimiento, protocolo de evaluaci{\'o}n y
  aspectos {\'e}ticos.
  \item \textbf{Cap{\'i}tulo~IV --- Factibilidad:} an{\'a}lisis de la factibilidad t{\'e}cnica,
  econ{\'o}mica y operativa.
  \item \textbf{Cap{\'i}tulo~V --- Aspectos administrativos:} cronograma, presupuesto y equipo.
\end{itemize}
\nuevo{\noindent La \textbf{matriz de consistencia}, incluida al cierre del Cap{\'i}tulo~III,
articula de forma integral el problema, los objetivos, las hip{\'o}tesis, las variables y la
metodolog{\'i}a.}
  % Cap. I  - Introduccion
% ##################################################################
% CAPÍTULO II: MARCO TEÓRICO Y ESTADO DEL ARTE
% ##################################################################
\clearpage
\section{Marco te{\'o}rico}

\subsection{Antecedentes}
Los trabajos previos sobre tecnolog{\'i}as del habla y traducci{\'o}n autom{\'a}tica para el quechua
se organizan en tres niveles ---internacional, nacional y local---. En conjunto evidencian avances
s{\'o}lidos por componente (ASR, MT y TTS), pero ninguno aborda su integraci{\'o}n en el dominio
m{\'e}dico-paciente ni la variante de Cusco (\sig{quz}).

\obsdos{A continuaci{\'o}n se desarrollan los antecedentes m{\'a}s relevantes siguiendo, en cada caso,
cinco elementos: (i) la cita y el t{\'i}tulo del estudio, (ii) su dise{\~n}o metodol{\'o}gico, (iii)
los instrumentos o modelos empleados, (iv) sus resultados y conclusi{\'o}n principal, y (v) su aporte
a la presente tesis. Al final, la Tabla~\ref{tab:sintesis} sintetiza el conjunto por componente.}

\subsubsection{Antecedentes internacionales}
\begin{bloqueObsDos}
\noindent\textbf{Antecedente 1.}
\begin{itemize}\setlength{\itemsep}{1pt}
  \item \textbf{Cita y t{\'i}tulo:} \textcite{ortega2020quechua}, \emph{Neural machine translation with a polysynthetic low resource language} (revista \emph{Machine Translation}); primer trabajo de traducci{\'o}n autom{\'a}tica neuronal Espa{\~n}ol--Quechua.
  \item \textbf{Dise{\~n}o metodol{\'o}gico:} Enfoque cuantitativo, dise{\~n}o experimental, orientado a superar la escasez de datos de una lengua polisint{\'e}tica.
  \item \textbf{Instrumentos:} Redes neuronales de traducci{\'o}n con segmentaci{\'o}n morfol{\'o}gica (subpalabras guiadas por la morfolog{\'i}a), evaluadas con la m{\'e}trica BLEU.
  \item \textbf{Resultados y conclusi{\'o}n:} La segmentaci{\'o}n morfol{\'o}gica mejora la traducci{\'o}n frente a la l{\'i}nea base; concluyen que el tratamiento expl{\'i}cito de la morfolog{\'i}a es decisivo en lenguas aglutinantes.
  \item \textbf{Aporte a la tesis:} Justifica ling{\"u}{\'i}sticamente el problema y fundamenta el tratamiento morfol{\'o}gico espec{\'i}fico del quechua al ajustar los modelos de traducci{\'o}n.
\end{itemize}

\noindent\textbf{Antecedente 2.}
\begin{itemize}\setlength{\itemsep}{1pt}
  \item \textbf{Cita y t{\'i}tulo:} \textcite{Ortega2023261,Ortega2025260}, serie \emph{QUESPA Submission for the IWSLT Dialectal and Low-resource Speech Translation Task} (IWSLT, 2023--2025); traducci{\'o}n de voz Quechua--Espa{\~n}ol.
  \item \textbf{Dise{\~n}o metodol{\'o}gico:} Experimental y comparativo entre ediciones sucesivas.
  \item \textbf{Instrumentos:} Arquitectura en cascada (reconocimiento de voz seguido de traducci{\'o}n) sobre modelos multiling{\"u}es preentrenados (MMS y NLLB) afinados por transferencia, evaluados con BLEU.
  \item \textbf{Resultados y conclusi{\'o}n:} El desempe{\~n}o subi{\'o} de cerca de 16 a \textbf{26.7 BLEU} (2023--2025), incorporando en 2025 la variante Collao/Cusco (\sig{quz}); concluyen que el ajuste de modelos preentrenados es la v{\'i}a m{\'a}s eficaz para esta lengua.
  \item \textbf{Aporte a la tesis:} Es el antecedente m{\'a}s cercano; aporta la arquitectura en cascada y el protocolo de evaluaci{\'o}n que se reutilizan, y revela la brecha a cerrar (la variante \sig{quz} y el dominio cl{\'i}nico, a{\'u}n no resueltos de extremo a extremo).
\end{itemize}
\end{bloqueObsDos}

\subsubsection{Antecedentes nacionales}
\begin{bloqueObsDos}
\noindent\textbf{Antecedente 3.}
\begin{itemize}\setlength{\itemsep}{1pt}
  \item \textbf{Cita y t{\'i}tulo:} \textcite{Zevallos202061}, \emph{Automatic Speech Recognition of Quechua Language Using HMM Toolkit} (Springer, CCIS); reconocedor de voz para el quechua.
  \item \textbf{Dise{\~n}o metodol{\'o}gico:} Cuantitativo-experimental.
  \item \textbf{Instrumentos:} Modelos ocultos de Markov (HMM/HTK) entrenados sobre un corpus de voz quechua, evaluados con la tasa de error de palabra (WER).
  \item \textbf{Resultados y conclusi{\'o}n:} \textbf{12.70\,\% de WER} y, con aumento de datos \parencite{Zevallos20223518}, 22.75\,\% de WER sobre datos m{\'a}s variados; concluyen que la escasez de datos es el principal limitante del ASR quechua.
  \item \textbf{Aporte a la tesis:} Establece una l{\'i}nea base nacional de WER para el m{\'o}dulo de reconocimiento y respalda el aumento de datos ante la escasez de corpus cl{\'i}nico.
\end{itemize}

\noindent\textbf{Antecedente 4.}
\begin{itemize}\setlength{\itemsep}{1pt}
  \item \textbf{Cita y t{\'i}tulo:} \textcite{HuamanCuro2026}, \emph{Quechua Speech Synthesizers} (Springer, CCIS); s{\'i}ntesis de voz para el quechua.
  \item \textbf{Dise{\~n}o metodol{\'o}gico:} Experimental y comparativo entre cuatro arquitecturas de TTS.
  \item \textbf{Instrumentos:} Modelos neuronales de s{\'i}ntesis entrenados sobre quechua chanca, valorados con el \emph{Mean Opinion Score} (MOS) mediante juicios perceptuales.
  \item \textbf{Resultados y conclusi{\'o}n:} MOS de hasta \textbf{4.9}, cercano a la voz humana; concluyen que es viable generar voz quechua de alta naturalidad con arquitecturas actuales.
  \item \textbf{Aporte a la tesis:} Confirma la factibilidad del m{\'o}dulo de salida por voz (TTS) ---clave para usuarios con baja alfabetizaci{\'o}n--- y valida el MOS como instrumento de medici{\'o}n.
\end{itemize}

\noindent\textbf{Antecedente 5.}
\begin{itemize}\setlength{\itemsep}{1pt}
  \item \textbf{Cita y t{\'i}tulo:} \textcite{CuevaMedina2024}, \emph{Development of a Neural Machine Translation Model Optimized with BERT for Translation from Quechua to Spanish} (LACCEI); traducci{\'o}n neuronal Quechua--Espa{\~n}ol.
  \item \textbf{Dise{\~n}o metodol{\'o}gico:} Aplicado, de dise{\~n}o experimental.
  \item \textbf{Instrumentos:} Modelo de traducci{\'o}n neuronal optimizado con representaciones tipo BERT, evaluado con m{\'e}tricas de traducci{\'o}n (BLEU/chrF).
  \item \textbf{Resultados y conclusi{\'o}n:} Mejoras al incorporar modelos preentrenados frente a enfoques previos; concluyen que la transferencia desde modelos de lenguaje mejora la traducci{\'o}n en contextos de bajo recurso.
  \item \textbf{Aporte a la tesis:} Confirma, en el contexto peruano, la estrategia de \emph{fine-tuning} sobre modelos preentrenados que aqu{\'i} se adopta para la adaptaci{\'o}n dialectal \sig{quy}$\to$\sig{quz}.
\end{itemize}
\end{bloqueObsDos}

\subsubsection{Antecedentes locales}
\begin{bloqueObsDos}
\noindent\textbf{Antecedente 6.}
\begin{itemize}\setlength{\itemsep}{1pt}
  \item \textbf{Cita y t{\'i}tulo:} \textcite{begazo2026}, \emph{Barreras ling{\"u}{\'i}sticas y culturales en la atenci{\'o}n hospitalaria: la experiencia de pacientes quechua-hablantes} (proyecto de investigaci{\'o}n); barrera idiom{\'a}tica en hospitales de referencia de Cusco.
  \item \textbf{Dise{\~n}o metodol{\'o}gico:} Estudio de campo centrado en la experiencia de los pacientes.
  \item \textbf{Instrumentos:} T{\'e}cnicas de recolecci{\'o}n sobre la interacci{\'o}n cl{\'i}nica real.
  \item \textbf{Resultados y conclusi{\'o}n:} Documenta c{\'o}mo la desconexi{\'o}n idiom{\'a}tica vulnera la relaci{\'o}n terap{\'e}utica y concluye que la barrera compromete la calidad y la seguridad de la atenci{\'o}n; en una regi{\'o}n andina vecina, \textcite{cuaila2022} hall{\'o} una relaci{\'o}n directa entre el dominio del quechua del personal y la calidad de atenci{\'o}n percibida.
  \item \textbf{Aporte a la tesis:} Aporta la evidencia local del problema que motiva la propuesta y respalda la pertinencia de una soluci{\'o}n por voz para el consultorio cusque{\~n}o.
\end{itemize}
\end{bloqueObsDos}

\noindent En el plano tecnol{\'o}gico local, el primer sistema de traducci{\'o}n neural para el
quechua se desarroll{\'o} con participaci{\'o}n de \emph{Hinantin Software} de Cusco
\parencite{ortega2020quechua}, y la variante Collao/Cusco (\sig{quz}) fue incorporada reci{\'e}n en
la campa{\~n}a QUESPA 2025 \parencite{Ortega2025260}. A ello se suma el reto de la alfabetizaci{\'o}n
en zonas rurales andinas \parencite{inei2023,dispurse2023}.

\obstres{Complementan este panorama otros
antecedentes nacionales e internacionales: la traducci{\'o}n de voz en cascada
Espa{\~n}ol$\to$Quechua \parencite{Alvarez-Crespo2023837}, los enfoques multiling{\"u}es para las lenguas
del Per{\'u} \parencite{Veliz2021241,Oncevay2021194}, la segmentaci{\'o}n morfol{\'o}gica
\parencite{Chen202120} y las representaciones autosupervisadas \parencite{Chen20246444}.}

\obstres{\noindent A ello se suman el aumento de datos
sint{\'e}ticos \parencite{Gow-Smith2023192,Ahmed2023156,Dhawan2026119,Ebrahimi2024236} y recursos
como \emph{QuBERT} \parencite{Zevallos20221}, \emph{Huqariq} \parencite{zevallos2022huqariq} y la API
\emph{Siminchikkunarayku} \parencite{9514251}, todos sintetizados en la Tabla~\ref{tab:sintesis}.}

No se identific{\'o} ning{\'u}n antecedente
local que integre reconocimiento, traducci{\'o}n y s{\'i}ntesis de voz para la atenci{\'o}n en salud,
lo que confirma la pertinencia de esta investigaci{\'o}n.

\subsection{Estado del arte y brecha de investigaci{\'o}n}
A partir de una \textbf{revisi{\'o}n sistem{\'a}tica} (Scopus, IEEE Xplore y ACM Digital Library;
62 registros identificados, 58 tras depurar duplicados y 30 evaluados a texto completo), los
antecedentes de la secci{\'o}n anterior se sintetizan por componente en la Tabla~\ref{tab:sintesis}
---sin repetir las citas ya detalladas---, lo que permite ubicar con precisi{\'o}n qu{\'e} se ha
resuelto y qu{\'e} sigue abierto.

\begin{table}[H]
\centering
\caption{S{\'i}ntesis del estado del arte por componente de la cascada.}
\label{tab:sintesis}
\small
\begin{tabularx}{\linewidth}{@{}l X l@{}}
\toprule\hdr{Componente} & \hdr{Avance / mejor resultado reportado} & \hdr{{\textquestiondown}Salud?} \\
\midrule
ASR (quechua)                  & WER 12.70\,\% (HMM/HTK) y 22.75\,\% con aumento de datos & No \\
MT (QU--ES)                    & Hasta $+$11\,\% chrF; BLEU 15.47 con modelos ajustados   & No \\
TTS (quechua)                  & MOS de hasta 4.9 (XTTS\,v2)                              & No \\
Traducci{\'o}n de voz (cascada) & 26.7 BLEU (QUESPA 2025, variante Collao)                & No \\
Recursos / corpus              & QuBERT, \emph{Huqariq}, API \emph{Siminchikkunarayku}    & Habilitador \\
\bottomrule
\end{tabularx}
\notaAPA{Elaboraci{\'o}n propia a partir de la revisi{\'o}n sistem{\'a}tica (Scopus, IEEE Xplore y ACM); s{\'i}ntesis por componente. La columna ``{\textquestiondown}Salud?'' indica si el trabajo abord{\'o} el dominio sanitario.}
\end{table}

La s{\'i}ntesis evidencia avances s{\'o}lidos por componente, pero ninguno en el \textbf{dominio
cl{\'i}nico} ni con la variante de Cusco (\sig{quz}) integrada de extremo a extremo. La
\textbf{brecha} es, por tanto, la intersecci{\'o}n ---hoy vac{\'i}a--- entre las tecnolog{\'i}as del
habla para el quechua y la comunicaci{\'o}n m{\'e}dico-paciente: no existe un corpus cl{\'i}nico
Quechua--Espa{\~n}ol ni una evaluaci{\'o}n emp{\'i}rica de estos modelos en salud, y NLLB a{\'u}n no
cubre \sig{quz}. Cerrar esa intersecci{\'o}n es el aporte de este proyecto.

\subsection{Fundamentos te{\'o}ricos}

\subsubsection{Tecnolog{\'i}as del habla y traducci{\'o}n autom{\'a}tica neuronal}
El \textbf{reconocimiento autom{\'a}tico de voz (ASR)} convierte una se{\~n}al de voz en texto; la
\textbf{traducci{\'o}n autom{\'a}tica (MT)}, hoy mayoritariamente \emph{neuronal} con arquitecturas
Transformer, traduce entre idiomas; y la \textbf{s{\'i}ntesis de voz (TTS)} genera voz a partir de
texto. En \textbf{lenguas de bajo recurso} ---como el quechua, de morfolog{\'i}a \glo{polisintetica}{polisint{\'e}tica}---
estas tareas se benefician de la transferencia desde modelos entrenados sobre grandes vol{\'u}menes de
datos multiling{\"u}es.

\begin{bloqueNuevo}
\noindent La base com{\'u}n de estos modelos es la arquitectura \emph{Transformer}
\parencite{vaswani2017attention}, cuyo n{\'u}cleo es el mecanismo de \emph{autoatenci{\'o}n}: cada
token se proyecta en consultas $Q$, claves $K$ y valores $V$, y la atenci{\'o}n pondera la
informaci{\'o}n relevante del contexto mediante
\begin{equation}
\mathrm{Attention}(Q,K,V) = \mathrm{softmax}\!\left(\frac{QK^{\top}}{\sqrt{d_k}}\right)V,
\label{eq:atencion}
\end{equation}
\regecuacion{Mecanismo de autoatenci{\'o}n del Transformer}
donde $d_k$ es la dimensi{\'o}n de las claves y el factor $1/\sqrt{d_k}$ estabiliza los gradientes.
Este mecanismo modela dependencias de largo alcance, decisivo en una lengua polisint{\'e}tica como el
quechua, donde una sola palabra concentra mucha informaci{\'o}n gramatical.
\end{bloqueNuevo}

\subsubsection{Modelos multiling{\"u}es preentrenados (MMS, NLLB, Whisper)}
\textbf{MMS} \parencite{pratap2024mms} cubre el quechua de Cusco (\sig{quz}) en reconocimiento y
s{\'i}ntesis; \textbf{NLLB} \parencite{costajussa2022nllb} ofrece traducci{\'o}n multiling{\"u}e que
incluye el quechua de Ayacucho (\sig{quy}); y \textbf{Whisper} \parencite{radford2023whisper} se
emplea como referencia de ASR con cobertura marginal del quechua. Estos modelos permiten reutilizar
conocimiento ling{\"u}{\'i}stico general y especializarlo mediante ajuste.

\nuevo{Conviene distinguir la \emph{teor{\'i}a} de la \emph{herramienta}: el Transformer
(ecuaci{\'o}n~\ref{eq:atencion}) es el modelo abstracto, mientras que MMS, NLLB y Whisper son
\emph{frameworks} preentrenados que lo implementan y que en esta tesis se reutilizan por ajuste
fino.}

\subsubsection{Ajuste fino y adaptaci{\'o}n dialectal (\texttt{quy}$\to$\texttt{quz})}
El \emph{ajuste fino} (\emph{fine-tuning}) adapta un modelo preentrenado a un dominio o variante
espec{\'i}ficos con relativamente pocos datos. El aporte t{\'e}cnico central de esta tesis es la
adaptaci{\'o}n del par de traducci{\'o}n del quechua de Ayacucho (\sig{quy}) al de Cusco
(\sig{quz}) sobre el vocabulario cl{\'i}nico, cerrando la brecha dialectal que hoy limita a NLLB.

\begin{bloqueNuevo}
\noindent Formalmente, partiendo de los pesos preentrenados $\theta_0$, el ajuste minimiza la
\emph{entrop{\'i}a cruzada} sobre el corpus de dominio $\mathcal{D}$:
\begin{equation}
\mathcal{L}(\theta) = -\!\!\sum_{(x,y)\in\mathcal{D}}\;\sum_{t=1}^{|y|}
\log P_{\theta}\!\left(y_t \mid y_{<t},\, x\right),
\label{eq:perdida}
\end{equation}
\regecuacion{Funci{\'o}n de p{\'e}rdida (entrop{\'i}a cruzada) del ajuste fino}
y los par{\'a}metros se actualizan por descenso de gradiente
$\theta \leftarrow \theta - \eta\,\nabla_{\theta}\mathcal{L}(\theta)$, con tasa de aprendizaje $\eta$.
De este modo, el ajuste reorienta la distribuci{\'o}n del modelo hacia el vocabulario cl{\'i}nico y la
variante \sig{quz}; esta es la base te{\'o}rica que explica la reducci{\'o}n esperada del WER y la
mejora de BLEU/chrF postuladas en la hip{\'o}tesis H2.
\end{bloqueNuevo}

\subsubsection{Arquitectura en cascada ASR$\to$MT$\to$TTS}
El encadenamiento de los tres componentes en una \textbf{arquitectura en cascada} habilita un
di{\'a}logo bidireccional por voz: la voz del hablante se transcribe (ASR), el texto se traduce (MT)
y la traducci{\'o}n se reproduce en la lengua del interlocutor (TTS). La modularidad facilita ajustar
y evaluar cada etapa por separado.

\begin{bloqueNuevo}
\noindent El flujo de datos es
$a \xrightarrow{\text{ASR}} t_s \xrightarrow{\text{MT}} t_t \xrightarrow{\text{TTS}} a_t$:
la se{\~n}al de audio $a$ se transcribe a texto origen $t_s$, se traduce a texto destino $t_t$ y se
sintetiza como audio $a_t$ (v{\'e}ase el Algoritmo~\ref{alg:pipeline} del Cap{\'i}tulo~III). Al ser
etapas secuenciales, la latencia por turno es aditiva,
$T = T_{\mathrm{ASR}} + T_{\mathrm{MT}} + T_{\mathrm{TTS}}$; el presupuesto objetivo se resume en la
Tabla~\ref{tab:latencias}, lo que justifica evaluar cada m{\'o}dulo por separado y su integraci{\'o}n.
\end{bloqueNuevo}

\begin{table}[H]
\centering
\caption{Presupuesto de latencia por etapa de la cascada (objetivo por turno).}
\label{tab:latencias}
\small
{
\begin{tabularx}{\linewidth}{@{}l X r@{}}
\toprule\hdr{Etapa} & \hdr{Operaci{\'o}n} & \hdr{Latencia objetivo (s)} \\
\midrule
ASR (MMS) & Voz $\to$ texto origen ($t_s$)        & $\le 1.5$ \\
MT (NLLB) & Texto origen $\to$ texto destino ($t_t$) & $\le 1.0$ \\
TTS (MMS) & Texto destino $\to$ voz ($a_t$)       & $\le 1.5$ \\
\multicolumn{2}{@{}l}{\textbf{Total por turno} ($T$)} & $\mathbf{\le 4.0}$ \\
\bottomrule
\end{tabularx}}
\notaAPA{Valores objetivo para el prototipo en equipo de gama media; se verifican emp{\'i}ricamente en el Cap{\'i}tulo~III (m{\'e}tricas de eficiencia).}
\end{table}

\subsubsection{M{\'e}tricas de evaluaci{\'o}n}
El desempe{\~n}o se mide con m{\'e}tricas estandarizadas: \textbf{WER} para el reconocimiento,
\textbf{BLEU} \parencite{papineni2002bleu} y \textbf{chrF} \parencite{popovic2015chrf} para la
traducci{\'o}n ---esta {\'u}ltima adecuada para lenguas morfol{\'o}gicamente ricas--- y \textbf{MOS}
para la naturalidad de la voz sintetizada.

\begin{bloqueNuevo}
\noindent Se definen formalmente a continuaci{\'o}n. La \emph{tasa de error de palabra} compara la
transcripci{\'o}n con la referencia:
\begin{equation}
\mathrm{WER} = \frac{S + D + I}{N},
\label{eq:wer}
\end{equation}
\regecuacion{Tasa de error de palabra (WER)}
con $S$, $D$ e $I$ sustituciones, eliminaciones e inserciones sobre $N$ palabras de referencia.
\emph{BLEU} mide la precisi{\'o}n de $n$-gramas con penalizaci{\'o}n por brevedad (BP):
\begin{equation}
\mathrm{BLEU} = \mathrm{BP}\cdot\exp\!\left(\textstyle\sum_{n=1}^{4} w_n \log p_n\right),
\label{eq:bleu}
\end{equation}
\regecuacion{Puntaje BLEU}
donde $p_n$ es la precisi{\'o}n de $n$-gramas y $w_n$ su peso. \emph{chrF}, m{\'a}s robusta en
morfolog{\'i}a rica, es la media arm{\'o}nica ponderada de precisi{\'o}n (chrP) y \emph{recall} (chrR)
de $n$-gramas de \emph{caracteres}:
\begin{equation}
\mathrm{chrF}_{\beta} = (1 + \beta^{2})\,\frac{\mathrm{chrP}\cdot\mathrm{chrR}}{\beta^{2}\,\mathrm{chrP} + \mathrm{chrR}}.
\label{eq:chrf}
\end{equation}
\regecuacion{Puntaje chrF ($F$-score de $n$-gramas de caracteres)}
Por {\'u}ltimo, el \emph{Mean Opinion Score} es el promedio de $J$ juicios perceptuales
$m_j \in [1,5]$, es decir $\mathrm{MOS} = \tfrac{1}{J}\sum_{j=1}^{J} m_j$. La aplicaci{\'o}n operativa
de estas m{\'e}tricas se retoma en el Cap{\'i}tulo~III.
\end{bloqueNuevo}

\subsection{Definici{\'o}n de t{\'e}rminos b{\'a}sicos}
\begin{itemize}
  \item \textbf{Corpus:} Conjunto estructurado de textos o grabaciones para entrenar o evaluar modelos.
  \item \textbf{Modelo preentrenado:} Modelo entrenado sobre grandes vol{\'u}menes de datos, reutilizable mediante ajuste.
  \item \textbf{Lengua de bajo recurso:} Idioma con escasez de datos digitales para el PLN.
  \item \textbf{Lengua polisint{\'e}tica:} Lengua que concentra mucha informaci{\'o}n gramatical en una sola palabra, como el quechua.
  \item \textbf{Pipeline:} Secuencia encadenada de etapas (voz $\to$ texto $\to$ traducci{\'o}n $\to$ voz).
\end{itemize}
       % Cap. II - Marco teorico
% ##################################################################
% CAPÍTULO III: METODOLOGÍA DE LA INVESTIGACIÓN
% ##################################################################
\clearpage
\section{Metodolog{\'i}a de la investigaci{\'o}n}

\subsection{Enfoque, tipo y dise{\~n}o de investigaci{\'o}n}

\subsubsection{Enfoque de investigaci{\'o}n}
La investigaci{\'o}n adopta un enfoque \textbf{\nuevo{cuantitativo}} \nuevo{\parencite{hernandez2018metodologia}:
todas las variables ---tanto las de desempe{\~n}o (WER, BLEU, chrF, MOS, latencia) como las de
percepci{\'o}n (escalas Likert y MOS)--- se miden con instrumentos num{\'e}ricos y se analizan con
estad{\'i}stica descriptiva e inferencial, sin recurrir a categor{\'i}as cualitativas.}

\subsubsection{Tipo y nivel de investigaci{\'o}n}
Siguiendo a \textcite{hernandez2018metodologia}, es de tipo \textbf{aplicado} ---resuelve un
problema concreto mediante el desarrollo de un prototipo, sin buscar teor{\'i}a nueva--- y de nivel
\textbf{aplicativo}, con alcance \emph{descriptivo} en la fase de diagn{\'o}stico (caracterizaci{\'o}n
de barreras y requerimientos) y \emph{explicativo} en la validaci{\'o}n, al contrastar el efecto del
sistema frente a la ausencia de apoyo tecnol{\'o}gico.

La investigaci{\'o}n se organiza en las tres
fases del modelo de ingenier{\'i}a ---\textbf{diagn{\'o}stico}, \textbf{construcci{\'o}n} y
\textbf{validaci{\'o}n}---, en correspondencia con los tres objetivos espec{\'i}ficos.

\subsubsection{Dise{\~n}o de investigaci{\'o}n}
El dise{\~n}o es \textbf{\obscuatro{experimental, de tipo cuasi-experimental con medidas repetidas
(intra-sujeto)}}: \obscuatro{se manipula la variable independiente ---uso del sistema y modalidad de
salida--- y cada unidad se mide en dos condiciones (con sistema \emph{vs.}\ sin sistema; texto
\emph{vs.}\ texto+voz), actuando como su propio control. Esta estructura de medidas pareadas
fortalece la comparaci{\'o}n causal frente a un dise{\~n}o pre-experimental de un solo grupo, sin
requerir la aleatorizaci{\'o}n de grupos ---inviable con una poblaci{\'o}n ling{\"u}{\'i}stica
acotada---}; la unidad de an{\'a}lisis es el \emph{turno de traducci{\'o}n} dentro de una
consulta simulada. La Figura~\ref{fig:arq} muestra el flujo del sistema y el
Algoritmo~\ref{alg:pipeline} detalla el procesamiento de un turno de di{\'a}logo.

\begin{figure}[H]
\centering
\caption{Pipeline modular \textbf{bidireccional} por voz; la traducci{\'o}n NLLB adaptada al quechua de Cusco (\sig{quz}) es el n{\'u}cleo del aporte.}
\label{fig:arq}
\begin{tikzpicture}[
    node distance=0.6cm and 0.7cm,
    boxA/.style={rectangle, rounded corners, draw=g4, thick, fill=g3!18, text width=2.0cm, align=center, minimum height=1.2cm, font=\scriptsize},
    core/.style={rectangle, rounded corners, draw=g8, very thick, fill=g8!22, text width=2.0cm, align=center, minimum height=1.2cm, font=\scriptsize},
    boxB/.style={rectangle, rounded corners, draw=g11, thick, fill=g12!22, text width=2.0cm, align=center, minimum height=1.2cm, font=\scriptsize},
    io/.style={ellipse, draw=g5, thick, fill=white, text width=1.8cm, align=center, font=\scriptsize},
    bi/.style={{Stealth[length=2.5mm]}-{Stealth[length=2.5mm]}, thick, draw=g5}
]
\node[io] (pac) {Paciente\\(voz quechua)};
\node[boxA, right=of pac] (v1) {MMS\\STT / TTS\\(\sig{quz})};
\node[core, right=of v1] (mt) {Traducci{\'o}n\\NLLB\\\sig{quz}$\leftrightarrow$es};
\node[boxB, right=of mt] (v2) {MMS\\STT / TTS\\(es)};
\node[io, right=of v2] (med) {M{\'e}dico\\(voz espa{\~n}ol)};
\draw[bi] (pac) -- (v1);
\draw[bi] (v1) -- (mt);
\draw[bi] (mt) -- (v2);
\draw[bi] (v2) -- (med);
\end{tikzpicture}
\notaAPA{Elaboraci{\'o}n propia. El bloque central (traducci{\'o}n NLLB \sig{quz}$\leftrightarrow$es) concentra la adaptaci{\'o}n dialectal, n{\'u}cleo del aporte de la tesis.}
\end{figure}

\begin{algorithm}[H]
\caption{Traducci{\'o}n bidireccional por voz de un turno de di{\'a}logo}
\label{alg:pipeline}
\begin{algorithmic}[1]
\Require Audio de entrada $a$; idioma origen $L_s$; idioma destino $L_t$
\Ensure Texto traducido $t_t$ y audio sintetizado $a_t$
\State $t_s \gets \textsc{STT}_{\mathrm{MMS}}(a,\,L_s)$
\State $t_t \gets \textsc{MT}_{\mathrm{NLLB}}(t_s,\,L_s \rightarrow L_t)$
\State Mostrar $t_t$ en pantalla
\State $a_t \gets \textsc{TTS}_{\mathrm{MMS}}(t_t,\,L_t)$
\State Reproducir $a_t$
\State \Return $(t_t,\,a_t)$
\end{algorithmic}
\end{algorithm}
\vspace{-9pt}
{\setlength{\parindent}{0pt}\notaAPA{Elaboraci{\'o}n propia, basada en la arquitectura en cascada ASR$\to$MT$\to$TTS descrita en el Cap{\'i}tulo~II.}}
\par\vspace{\baselineskip}

\begin{bloqueNuevo}
\noindent Las funciones $\textsc{STT}$, $\textsc{MT}$ y $\textsc{TTS}$ no son cajas negras: internamente
aplican la autoatenci{\'o}n del Transformer (ecuaci{\'o}n~\ref{eq:atencion}) sobre secuencias de
\emph{tokens}, con los pesos $\theta$ previamente especializados por \emph{fine-tuning}
(ecuaci{\'o}n~\ref{eq:perdida}). As{\'i}, el paso~2 (MT con NLLB) es el que concentra la
adaptaci{\'o}n dialectal \sig{quy}$\to$\sig{quz}, aporte central de la tesis.
\end{bloqueNuevo}

\subsection{\obsdos{Hip{\'o}tesis}}
\begin{bloqueObsDos}
\noindent\textbf{Hip{\'o}tesis general:} Un sistema bidireccional de traducci{\'o}n de voz
Quechua--Espa{\~n}ol, basado en aprendizaje profundo (modelos preentrenados basados en
Transformer), mejora significativamente la comunicaci{\'o}n m{\'e}dico-paciente durante las consultas
m{\'e}dicas en Cusco frente a la ausencia de apoyo tecnol{\'o}gico.
\begin{itemize}
  \item \textbf{H1 (Diagn{\'o}stico} ---\obscuatro{proposici{\'o}n descriptiva}\textbf{):} \obscuatro{La
  comunicaci{\'o}n m{\'e}dico-paciente con pacientes quechuahablantes presenta barreras
  ling{\"u}{\'i}sticas \emph{caracterizables}, cuya identificaci{\'o}n define los requisitos ---l{\'e}xico
  cl{\'i}nico y escenarios de triaje, s{\'i}ntomas e indicaciones--- del corpus Quechua--Espa{\~n}ol.
  Al ser diagn{\'o}stica, H1 se verifica por caracterizaci{\'o}n (cobertura l{\'e}xica y concordancia
  entre expertos), no por un contraste estad{\'i}stico de medias.}
  \item \textbf{H2 (Construcci{\'o}n):} El ajuste (\emph{fine-tuning}) de los modelos preentrenados integrados en el \emph{pipeline} reduce el WER y mejora el BLEU/chrF frente a los modelos sin ajustar.
  \item \textbf{H3 (Validaci{\'o}n):} El uso del sistema mejora significativamente la precisi{\'o}n, la calidad de la voz y la comprensi{\'o}n mutua en la comunicaci{\'o}n m{\'e}dico-paciente frente a la ausencia de apoyo tecnol{\'o}gico, sobre todo con baja alfabetizaci{\'o}n.
\end{itemize}
\noindent \obscuatro{H2 y H3 son hip{\'o}tesis \emph{contrastables} ---se someten a prueba estad{\'i}stica
(\emph{t} de Student o Wilcoxon, seg{\'u}n normalidad)---, mientras que H1 es una proposici{\'o}n
diagn{\'o}stica que se valida por caracterizaci{\'o}n.} Todas se articulan de forma integral en la
matriz de consistencia (Tabla~\ref{tab:matriz-consist}).
\end{bloqueObsDos}

\subsection{Variables y operacionalizaci{\'o}n}

\subsubsection{Identificaci{\'o}n de variables}
La \textbf{variable independiente} es \obsdos{el sistema de traducci{\'o}n (con o sin sistema)} y la
modalidad de salida (texto / texto+voz). Las \textbf{variables dependientes} son la precisi{\'o}n de
transcripci{\'o}n y traducci{\'o}n, la calidad de la voz y la comprensi{\'o}n/percepci{\'o}n. El
nivel de alfabetizaci{\'o}n act{\'u}a como \textbf{variable interviniente}.

\subsubsection{Matriz de operacionalizaci{\'o}n de variables}
La Tabla~\ref{tab:variables} resume la operacionalizaci{\'o}n.

\begin{table}[H]
\centering
\caption{Operacionalizaci{\'o}n de variables.}
\label{tab:variables}
\small
\begin{tabularx}{\linewidth}{@{}l X l@{}}
\toprule\hdr{Variable} & \hdr{Definici{\'o}n / dimensi{\'o}n} & \hdr{Indicador} \\
\midrule
\multirow{2}{*}{Independiente} & Uso del sistema de traducci{\'o}n & Con / sin sistema \\
                               & Modalidad de salida & Texto / texto+voz \\
\multirow{3}{*}{Dependiente}   & Precisi{\'o}n de transcripci{\'o}n & WER \\
                               & Precisi{\'o}n de traducci{\'o}n & BLEU / chrF \\
                               & Comprensi{\'o}n y percepci{\'o}n & \% aciertos, Likert, MOS \\
Interviniente                  & Nivel de alfabetizaci{\'o}n & Alto / bajo \\
\bottomrule
\end{tabularx}
\notaAPA{Elaboraci{\'o}n propia. WER: tasa de error de palabra; BLEU/chrF: m{\'e}tricas de traducci{\'o}n; MOS: \emph{Mean Opinion Score} (ecuaciones del Cap{\'i}tulo~II).}
\end{table}

\subsection{Unidades de an{\'a}lisis y selecci{\'o}n de consultas}

\subsubsection{Unidad de an{\'a}lisis}
La unidad de an{\'a}lisis es el \emph{turno de traducci{\'o}n} dentro de una consulta simulada
(una intervenci{\'o}n hablada que el sistema reconoce, traduce y reproduce). Los escenarios cubren
triaje, s{\'i}ntomas e indicaciones.

\subsubsection{Poblaci{\'o}n y muestra}
\begin{bloqueObsDos}
\noindent\textbf{Poblaci{\'o}n.} Es el universo de enunciados o \emph{turnos} de habla propios de la
comunicaci{\'o}n m{\'e}dico-paciente en consultas generales ---triaje, s{\'i}ntomas e
indicaciones--- en el par quechua de Cusco (\sig{quz})--espa{\~n}ol. Al tratarse de un universo
\emph{ling{\"u}{\'i}stico} y no de un conjunto de personas, se cuantifica en unidades de di{\'a}logo
(turnos), no en individuos.

\noindent\textbf{Muestra.} Es el \emph{corpus cl{\'i}nico} construido \emph{ad hoc} para el estudio:
un banco de frases-gui{\'o}n por escenario con sus grabaciones de audio, del orden de 30--50
frases-gui{\'o}n por escenario (triaje, s{\'i}ntomas e indicaciones) ---aproximadamente 90--150 en
total---, revisadas y validadas con hablantes nativos de la variante de Cusco.

\noindent\obscuatro{\textbf{Partici{\'o}n y control de fuga de datos.} El corpus se divide en los
subconjuntos de \emph{entrenamiento} (70\,\%), \emph{validaci{\'o}n} (15\,\%) y \emph{prueba}
(15\,\%), con separaci{\'o}n estricta: el conjunto de prueba \emph{no} participa en el ajuste, de modo
que las m{\'e}tricas (WER, BLEU, chrF) se reportan siempre sobre datos no vistos y no quedan infladas
por sobreajuste. Dado el tama{\~n}o reducido, la estimaci{\'o}n se robustece con \emph{validaci{\'o}n
cruzada} $k$-\emph{fold} ($k=5$), reportando la media y la desviaci{\'o}n de cada m{\'e}trica.}

\noindent\obscuatro{\textbf{Rol del corpus y aumento de datos.} Las 90--150 frases-gui{\'o}n
constituyen, ante todo, un \emph{banco de evaluaci{\'o}n} (\emph{benchmark}) de dominio cl{\'i}nico;
por s{\'i} solas son insuficientes como {\'u}nico material de entrenamiento para el \emph{fine-tuning}
de NLLB. Por ello el conjunto de ajuste se amplifica con \emph{aumento de datos} ---retrotraducci{\'o}n
(\emph{back-translation}) y generaci{\'o}n de variantes sint{\'e}ticas--- y se complementa con recursos
existentes para el quechua (\emph{Huqariq}, QuBERT), reservando el \emph{benchmark} cl{\'i}nico para la
medici{\'o}n final.}

\noindent\textbf{Unidad muestral (registro).} Cada \emph{turno} de traducci{\'o}n constituye una fila
del conjunto de datos, con sus campos de evaluaci{\'o}n: direcci{\'o}n \sig{quz}$\leftrightarrow$es,
transcripci{\'o}n, traducci{\'o}n de referencia, WER, BLEU, chrF, MOS y latencia (v{\'e}ase el
Anexo~1).

\noindent\textbf{Tipo de muestreo.} No probabil{\'i}stico por conveniencia, apropiado para una lengua
de bajo recurso sin corpus cl{\'i}nico preexistente: las frases se seleccionan por su pertinencia al
dominio y por su cobertura de los escenarios de consulta, no al azar.

\noindent\textbf{Participantes.} Para la validaci{\'o}n de percepci{\'o}n (comprensi{\'o}n y calidad
de voz) se convoca a 10--15 voluntarios biling{\"u}es \sig{quz}--espa{\~n}ol ---estudiantes de
ciencias de la salud o personal asistencial--- que act{\'u}an los roles en las consultas simuladas.
\emph{No} se emplean pacientes ni datos cl{\'i}nicos reales, en concordancia con los aspectos
{\'e}ticos del estudio.

\noindent\obscuatro{\textbf{Justificaci{\'o}n del tama{\~n}o.} El n{\'u}mero de jueces (10--15)
satisface el m{\'i}nimo habitual para estimar el MOS de forma estable ($\ge 10$ evaluadores); el
n{\'u}mero de turnos (90--150) responde a la disponibilidad real de hablantes nativos de la variante
de Cusco y a la naturaleza acotada del dominio. Al tratarse de una muestra peque{\~n}a, los resultados
se acompa{\~n}an de intervalos de confianza y del tama{\~n}o del efecto (\emph{d} de Cohen), y el
dise{\~n}o de medidas repetidas ---cada unidad como su propio control--- incrementa la potencia
estad{\'i}stica frente a un contraste entre grupos independientes.}
\end{bloqueObsDos}

\subsubsection{\nuevo{Entorno tecnol{\'o}gico (\emph{setup} experimental)}}
\begin{bloqueNuevo}
El ajuste y la evaluaci{\'o}n se ejecutan en el entorno de la Tabla~\ref{tab:setup}. El
\emph{fine-tuning} de MMS y NLLB requiere GPU; la inferencia del \emph{pipeline} integrado se valida
adem{\'a}s en equipo de gama media, compatible con un piloto en el punto de atenci{\'o}n.
\end{bloqueNuevo}

\begin{bloqueObsCuatro}
\noindent El \emph{fine-tuning} ---fase intensiva en c{\'o}mputo--- se ejecuta en \textbf{GPU en la
nube} (Google Colab Pro, Kaggle o cr{\'e}ditos de nube institucional), ya que el equipo local de
desarrollo (Windows\,11 con WSL2, sin GPU dedicada) solo soporta con holgura la \emph{inferencia} del
\emph{pipeline} integrado. Para acotar el costo, el ajuste emplea t{\'e}cnicas de bajo consumo de
memoria (\emph{LoRA}/\emph{PEFT} y precisi{\'o}n reducida), de modo que cada modelo se afina dentro de
los l{\'i}mites de VRAM de una GPU de nivel medio (p.\,ej.\ T4 de 16\,GB); el presupuesto de
horas-GPU se detalla en el Cap{\'i}tulo~V.
\end{bloqueObsCuatro}

\begin{table}[H]
\centering
\caption{\nuevo{Entorno de hardware y software del experimento.}}
\label{tab:setup}
\small
{
\begin{tabularx}{\linewidth}{@{}l X@{}}
\toprule\hdr{Componente} & \hdr{Especificaci{\'o}n prevista} \\
\midrule
GPU            & NVIDIA con $\ge 12$\,GB VRAM (p.\,ej.\ RTX 3060/T4); cr{\'e}ditos en la nube para fases intensivas \\
CPU / RAM      & 8 n{\'u}cleos / $\ge 16$\,GB \\
Almacenamiento & $\ge 256$\,GB SSD \\
Sistema        & Linux (Ubuntu 22.04) o Windows 11 con WSL2 \\
Lenguaje       & Python 3.11 \\
\emph{Frameworks} & PyTorch 2.x, Hugging Face \texttt{transformers}, \texttt{datasets}, \texttt{peft} (LoRA) \\
Evaluaci{\'o}n & \texttt{jiwer} (WER), \texttt{sacrebleu} (BLEU/chrF); MOS por panel \\
\bottomrule
\end{tabularx}}
\notaAPA{Las herramientas corresponden a los modelos justificados en el Cap{\'i}tulo~II (MMS, NLLB, Whisper) y a sus ecosistemas de c{\'o}digo abierto.}
\end{table}

\subsection{T{\'e}cnicas e instrumentos de recolecci{\'o}n}
\begin{itemize}
  \item \textbf{Corpus:} Lista de vocabulario y frases de triaje/s{\'i}ntomas/indicaciones, construida y revisada con personal de salud (fichas de validaci{\'o}n de contenido).
  \item \textbf{Registro de desempe{\~n}o:} \emph{Logs} autom{\'a}ticos del sistema (WER, BLEU, latencia por turno).
  \item \textbf{Encuesta de percepci{\'o}n:} Cuestionario tipo Likert (utilidad, confianza, comprensi{\'o}n) al personal y a los pacientes-actores.
  \item \textbf{Prueba de comprensi{\'o}n:} Tarea de recuperaci{\'o}n de la indicaci{\'o}n (\% de aciertos) para contrastar H2 y H3.
\end{itemize}

\subsection{Procedimiento metodol{\'o}gico}
\begin{enumerate}
  \item Construcci{\'o}n y validaci{\'o}n del corpus m{\'e}dico Quechua--Espa{\~n}ol.
  \item Ajuste (\emph{fine-tuning}) de MMS (reconocimiento y s{\'i}ntesis) y, como aporte central, de NLLB para el par quechua de Cusco (\sig{quz})--Espa{\~n}ol.
  \item Integraci{\'o}n del \emph{pipeline} y de la salida diferenciada por voz y texto.
  \item Ejecuci{\'o}n de consultas simuladas bajo las condiciones experimentales.
  \item Registro de m{\'e}tricas y aplicaci{\'o}n de instrumentos de percepci{\'o}n.
\end{enumerate}

\subsection{Protocolo de evaluaci{\'o}n}

\subsubsection{M{\'e}tricas de reconocimiento y traducci{\'o}n}
La precisi{\'o}n del reconocimiento se mide con la \textbf{tasa de error de palabra (WER)} y la de la
traducci{\'o}n con \textbf{BLEU} \parencite{papineni2002bleu} y \textbf{chrF} \parencite{popovic2015chrf},
adecuada para lenguas morfol{\'o}gicamente ricas. \nuevo{Sus formulaciones se definieron en el
Cap{\'i}tulo~II (ecuaciones~\ref{eq:wer}, \ref{eq:bleu} y~\ref{eq:chrf}).} Operativamente, el WER se
calcula por turno sobre la transcripci{\'o}n del ASR y BLEU/chrF sobre la traducci{\'o}n del MT, en
ambos sentidos \sig{quz}$\leftrightarrow$es, comparando contra referencias validadas por hablantes
nativos.

\subsubsection{M{\'e}tricas de calidad de voz y comprensi{\'o}n}
La calidad de la voz sintetizada se valora con el \emph{Mean Opinion Score} (MOS). La comprensi{\'o}n
se mide con el porcentaje de aciertos en la recuperaci{\'o}n de la indicaci{\'o}n y con escalas Likert
de utilidad y confianza.

\subsubsection{M{\'e}tricas de eficiencia}
La eficiencia se cuantifica con la \emph{latencia por turno} (tiempo entre el fin de la voz de
entrada y el inicio de la voz de salida) y el consumo de recursos de c{\'o}mputo.

\subsubsection{\obscuatro{L{\'i}neas base (\emph{baselines}) de comparaci{\'o}n}}
\begin{bloqueObsCuatro}
\noindent El desempe{\~n}o del sistema ajustado se interpreta siempre contra referencias expl{\'i}citas,
no en abstracto. Se contrastan tres condiciones: (i)~\textbf{modelos preentrenados sin ajuste}
---NLLB con el quechua de Ayacucho (\sig{quy}) y MMS base---, que fijan el punto de partida t{\'e}cnico;
(ii)~el \textbf{sistema ajustado} a la variante de Cusco (\sig{quz}) y al dominio cl{\'i}nico, que es
la propuesta; y (iii)~la \textbf{ausencia de apoyo tecnol{\'o}gico}, como referencia de la
comprensi{\'o}n m{\'e}dico-paciente. La mejora de (ii) sobre (i) contrasta H2 y la de (ii) sobre (iii)
contrasta H3.
\end{bloqueObsCuatro}

\subsubsection{\obscuatro{Criterios de {\'e}xito (metas cuantitativas)}}
\begin{bloqueObsCuatro}
\noindent El {\'e}xito del prototipo se declara verificable mediante umbrales objetivo, fijados a
partir de los resultados reportados en los antecedentes (Cap{\'i}tulo~II) y ajustados al dominio
cl{\'i}nico y a la variante de Cusco (Tabla~\ref{tab:metas}).
\end{bloqueObsCuatro}

\begin{table}[H]
\centering
\caption{\obscuatro{Metas de desempe{\~n}o del prototipo (criterios de {\'e}xito).}}
\label{tab:metas}
\small
\begin{tabularx}{\linewidth}{@{}l X l@{}}
\toprule\hdr{M{\'e}trica} & \hdr{Componente / sentido} & \hdr{Meta objetivo} \\
\midrule
WER            & Reconocimiento (ASR), menor es mejor      & $\le 25$\,\% \\
BLEU           & Traducci{\'o}n (MT), mayor es mejor         & $\ge 20$ \\
chrF           & Traducci{\'o}n (MT), morfolog{\'i}a rica    & $\ge 45$ \\
MOS            & Calidad de voz (TTS), escala 1--5         & $\ge 3.5$ \\
Comprensi{\'o}n & \% de aciertos en la indicaci{\'o}n        & $\ge 80$\,\% \\
Latencia       & Tiempo por turno (cascada completa)       & $\le 4.0$\,s \\
\bottomrule
\end{tabularx}
\notaAPA{\obscuatro{Elaboraci{\'o}n propia. Umbrales de referencia para declarar viable el prototipo;
se contrastan con las l{\'i}neas base y pueden reajustarse tras la fase de diagn{\'o}stico. Valores
anclados en los antecedentes del Cap{\'i}tulo~II (p.\,ej.\ 26.7 BLEU en QUESPA\,2025; MOS hasta 4.9 en
s{\'i}ntesis quechua).}}
\end{table}

\subsection{M{\'e}todos de an{\'a}lisis de datos}
Estad{\'i}stica descriptiva para WER/BLEU/MOS/latencia y pruebas de comparaci{\'o}n de
medias/proporciones (t de Student o Wilcoxon seg{\'u}n normalidad), con un nivel de significancia
$\alpha = 0{,}05$\label{sym:alpha}, para contrastar las hip{\'o}tesis (modelo ajustado vs. sin ajustar; texto vs.
texto+voz).

\nuevo{La elecci{\'o}n entre ambas pruebas se decide con un contraste previo de
normalidad (Shapiro--Wilk): si los datos son normales se aplica la \emph{t} de Student
(param{\'e}trica); si no, la prueba de Wilcoxon (no param{\'e}trica, basada en rangos), adecuada para
muestras peque{\~n}as como la de este estudio.}

\subsection{Aspectos {\'e}ticos y de gesti{\'o}n de datos}
El estudio usa {\'u}nicamente di{\'a}logos \textbf{simulados/actuados}, sin informaci{\'o}n cl{\'i}nica
real, evitando la necesidad de aprobaci{\'o}n de un comit{\'e} de {\'e}tica hospitalario para datos de
pacientes. Se obtiene consentimiento informado de los participantes voluntarios y se anonimizan las
grabaciones, en cumplimiento de los principios {\'e}ticos de la investigaci{\'o}n. Los datos se
almacenan de forma segura y se usan solo con fines de investigaci{\'o}n.

\begin{bloqueObsCuatro}
\noindent Aunque no se manejan datos de pacientes, el proyecto se acoge a la categor{\'i}a de
\textbf{riesgo m{\'i}nimo} y tramita la conformidad de la instancia correspondiente de la UNSAAC
---la Escuela Profesional de Ingenier{\'i}a Inform{\'a}tica y de Sistemas y, de requerirse, el
comit{\'e} institucional de {\'e}tica de la investigaci{\'o}n--- conforme al reglamento de grados y
t{\'i}tulos vigente. El consentimiento informado (Anexo~4), la participaci{\'o}n voluntaria con
derecho de retiro sin perjuicio y la anonimizaci{\'o}n de las grabaciones constituyen las garant{\'i}as
{\'e}ticas del estudio.
\end{bloqueObsCuatro}

\begin{landscape}
\subsection{\nuevo{Matriz de consistencia}}
\begin{bloqueNuevo}
\noindent Cada fila articula, en horizontal, problema, objetivo, hip{\'o}tesis, variables, indicadores y m{\'e}todo de an{\'a}lisis (Tabla~\ref{tab:matriz-consist}).
\end{bloqueNuevo}
\vspace{1pt}

\captionof{table}{\nuevo{Matriz de consistencia del proyecto.}}
\label{tab:matriz-consist}

% ----- barra de Tema (APA: sin relleno de color, solo borde) -----
\setlength{\fboxsep}{6pt}%
\noindent\fcolorbox{black}{white}{\parbox{\dimexpr\linewidth-2\fboxsep-2\fboxrule\relax}{\centering
  \footnotesize\textbf{Tema:} Desarrollo de un sistema bidireccional de traducci{\'o}n de voz Quechua--Espa{\~n}ol mediante aprendizaje profundo para consultas m{\'e}dicas en Cusco}}
\vspace{2pt}

\fontsize{7}{8.3}\selectfont
\renewcommand{\arraystretch}{1.08}
\setlength{\tabcolsep}{3.5pt}
\arrayrulecolor{black}
\noindent
\begin{tabularx}{\linewidth}{@{}
  >{\centering\arraybackslash}m{2.1cm}
  *{6}{>{\raggedright\arraybackslash}X}@{}}
\toprule
\hdr{Nivel} & \hdr{Problema} & \hdr{Objetivo} & \hdr{Hip{\'o}tesis} & \hdr{Variables} & \hdr{Dimensiones / indicadores} & \hdr{Metodolog{\'i}a} \\
\midrule
\celnivel{General} &
{\textquestiondown}C{\'o}mo desarrollar el sistema bidireccional de traducci{\'o}n de voz QU--ES para la comunicaci{\'o}n m{\'e}dico-paciente, mediante aprendizaje profundo, en Cusco? &
Desarrollar el sistema mediante el ajuste de modelos preentrenados (MMS, NLLB) para la comunicaci{\'o}n m{\'e}dico-paciente en Cusco. &
Mejora significativamente la comunicaci{\'o}n m{\'e}dico-paciente frente a la ausencia de apoyo. &
\textbf{VI:} Uso y modalidad (texto / texto+voz). \textbf{VD:} Precisi{\'o}n, voz, comprensi{\'o}n. &
WER, BLEU, chrF, MOS; \% comprensi{\'o}n; latencia (s). &
Aplicada; cuantitativo; \obscuatro{cuasi-exp.\ (medidas repetidas)}; t / Wilcoxon ($\alpha=0{,}05$). \\
\celnivel{Espec{\'i}fico 1 (Diagn{\'o}stico)} &
{\textquestiondown}Cu{\'a}les son las barreras de comunicaci{\'o}n y los requisitos del corpus cl{\'i}nico QU--ES? &
Caracterizar las barreras y los requerimientos del corpus cl{\'i}nico QU--ES. &
\obscuatro{(Descriptiva)} Las barreras definen los requisitos del corpus cl{\'i}nico para el ajuste. &
Barreras; corpus (calidad y cobertura). &
Cobertura l{\'e}xica; \% de concordancia. &
Nivel descriptivo; juicio de expertos; fichas de validaci{\'o}n. \\
\celnivel{Espec{\'i}fico 2 (Construcci{\'o}n)} &
{\textquestiondown}C{\'o}mo desarrollar el sistema (corpus + \emph{fine-tuning} de MMS y NLLB \sig{quy}$\to$\sig{quz} + \emph{pipeline})? &
Desarrollar el sistema: corpus + \emph{fine-tuning} de MMS y NLLB (\sig{quy}$\to$\sig{quz}), integrados en el \emph{pipeline}. &
El \emph{fine-tuning} reduce el WER y mejora el BLEU/chrF frente a los modelos sin ajustar. &
Modelo (base vs.\ ajustado); \emph{pipeline}. &
WER; BLEU; chrF; latencia (s). &
Nivel aplicativo; \emph{logs} autom{\'a}ticos; t / Wilcoxon (ajustado vs.\ base). \\
\celnivel{Espec{\'i}fico 3 (Validaci{\'o}n)} &
{\textquestiondown}En qu{\'e} medida el sistema mejora la precisi{\'o}n, la voz y la comprensi{\'o}n frente a la ausencia de apoyo? &
Evaluar la mejora en precisi{\'o}n, voz y comprensi{\'o}n en consultas simuladas. &
Mejora significativamente la precisi{\'o}n, la voz y la comprensi{\'o}n mutua. &
Precisi{\'o}n; voz; comprensi{\'o}n; percepci{\'o}n. &
WER/BLEU/chrF; MOS; \% de aciertos; Likert. &
Nivel explicativo; experimental; estad{\'i}stica inferencial. \\
\bottomrule
\end{tabularx}
\arrayrulecolor{black}
\notaAPA{Elaboraci{\'o}n propia. Las tres filas espec{\'i}ficas siguen el modelo de ingenier{\'i}a de tres fases (diagn{\'o}stico, construcci{\'o}n y validaci{\'o}n), en correspondencia 1 a 1 con los objetivos. VI/VD: variable independiente/dependiente. Las f{\'o}rmulas de WER, BLEU, chrF y MOS se definen en el Cap{\'i}tulo~II.}
\end{landscape}
        % Cap. III- Metodologia
% ##################################################################
% CAPÍTULO IV: ESTUDIO DE FACTIBILIDAD
% ##################################################################
\clearpage
\section{Estudio de factibilidad}

\subsection{Enfoque adoptado}
El proyecto es viable porque reutiliza modelos preentrenados de c{\'o}digo abierto (MMS, NLLB,
Whisper) mediante ajuste fino, lo que reduce dr{\'a}sticamente los requisitos de datos y c{\'o}mputo
frente a entrenar desde cero. La factibilidad se analiza en sus dimensiones t{\'e}cnica,
econ{\'o}mica y operativa.

\subsection{Factibilidad t{\'e}cnica}

\subsubsection{Hardware, c{\'o}mputo e inferencia}
El ajuste de modelos preentrenados puede realizarse en una estaci{\'o}n con GPU y complementarse con
cr{\'e}ditos de GPU en la nube para las fases intensivas. La inferencia del \emph{pipeline} integrado
es ejecutable en equipo de gama media, compatible con un piloto en el punto de atenci{\'o}n.
\obscuatro{En concreto, el \emph{fine-tuning} se realiza en \textbf{GPU en la nube} (Colab Pro,
Kaggle o nube institucional) con t{\'e}cnicas de bajo consumo de memoria (\emph{LoRA}/\emph{PEFT}),
dado que el equipo local (WSL2 sin GPU dedicada) se reserva para la inferencia; el presupuesto de
horas-GPU se refleja en la factibilidad econ{\'o}mica.}

\subsubsection{Estrategia de datos}
El cuello de botella es el corpus cl{\'i}nico Quechua--Espa{\~n}ol, que se construir{\'a} de forma
acotada (vocabulario y frases de triaje, s{\'i}ntomas e indicaciones) y se complementar{\'a} con
recursos existentes (\emph{Huqariq}, QuBERT) y datos sint{\'e}ticos cuando sea pertinente.

\subsubsection{Herramientas y software}
Se emplear{\'a}n bibliotecas de c{\'o}digo abierto para habla y traducci{\'o}n (modelos MMS, NLLB y
Whisper de Meta/OpenAI y sus ecosistemas), herramientas de evaluaci{\'o}n (c{\'a}lculo de WER, BLEU y
chrF) y entornos reproducibles de \emph{fine-tuning}.

\subsection{Factibilidad econ{\'o}mica}
El presupuesto estimado (Tabla~\ref{tab:presupuesto}) asciende a \textbf{S/ 7\,350.00}, dominado por
hardware con GPU y cr{\'e}ditos de c{\'o}mputo; es asumible en el marco de un trabajo de tesis y
puede financiarse con recursos propios y/o apoyo institucional.

\subsection{Factibilidad operativa}
La validaci{\'o}n con consultas simuladas y participantes voluntarios evita la dependencia de
aprobaciones cl{\'i}nicas complejas y permite ejecutar los experimentos \obscuatro{dentro del
cronograma anual (12 meses)}, en un centro de salud de referencia en Cusco.

\subsection{\obscuatro{An{\'a}lisis de riesgos y mitigaci{\'o}n}}
\begin{bloqueObsCuatro}
\noindent Los principales riesgos del proyecto y sus medidas de mitigaci{\'o}n se resumen en la
Tabla~\ref{tab:riesgos}; se priorizan por su impacto sobre los objetivos espec{\'i}ficos.
\end{bloqueObsCuatro}

\begin{table}[H]
\centering
\caption{\obscuatro{Matriz de riesgos y mitigaci{\'o}n.}}
\label{tab:riesgos}
\small
\begin{tabularx}{\linewidth}{@{}p{3.5cm} c X@{}}
\toprule\hdr{Riesgo} & \hdr{Prob./Imp.} & \hdr{Mitigaci{\'o}n} \\
\midrule
Corpus cl{\'i}nico insuficiente para el ajuste & Alta / Alto &
Aumento de datos (retrotraducci{\'o}n, sint{\'e}ticos) y reuso de \emph{Huqariq}/QuBERT; el corpus se usa sobre todo como \emph{benchmark}. \\
Escasa disponibilidad de hablantes nativos \sig{quz} & Media / Alto &
Convocatoria temprana de voluntarios biling{\"u}es; validaci{\'o}n por paneles reducidos ($\ge$ 3 revisores). \\
Latencia del \emph{pipeline} superior a la meta & Media / Medio &
Modelos ligeros, cuantizaci{\'o}n y evaluaci{\'o}n por etapas (presupuesto de la Tabla~\ref{tab:latencias}). \\
Costo/acceso a GPU en la nube & Media / Medio &
\emph{LoRA}/\emph{PEFT} y precisi{\'o}n reducida; cr{\'e}ditos gratuitos (Colab/Kaggle) y nube institucional. \\
Baja ganancia dialectal \sig{quy}$\to$\sig{quz} & Media / Alto &
Segmentaci{\'o}n morfol{\'o}gica, validaci{\'o}n con nativos y l{\'i}neas base de comparaci{\'o}n. \\
\bottomrule
\end{tabularx}
\notaAPA{\obscuatro{Elaboraci{\'o}n propia. Prob./Imp.: probabilidad e impacto estimados.}}
\end{table}

\subsection{Viabilidad y sostenibilidad}
El proceso (corpus + ajuste + evaluaci{\'o}n) es replicable en otras variantes del quechua o en otros
dominios, y el uso de modelos de c{\'o}digo abierto favorece la continuidad y la apropiaci{\'o}n local
de la tecnolog{\'i}a.
       % Cap. IV - Factibilidad
% ##################################################################
% CAPÍTULO V: ASPECTOS ADMINISTRATIVOS
% ##################################################################
\clearpage
\section{Aspectos administrativos}

\subsection{Cronograma de actividades}
El proyecto se ejecuta en \textbf{12 meses (octubre 2026 -- septiembre 2027)}, organizado por meses.
La Tabla~\ref{tab:cronograma} presenta el diagrama de Gantt, con las actividades agrupadas en tres
fases ---\emph{administrativa}, de \emph{desarrollo} (ASR, MT y TTS) e \emph{investigaci{\'o}n y
cierre}---; los m{\'o}dulos ASR, MT y TTS se ejecutan de forma parcialmente paralela mediante el
ajuste de modelos preentrenados.

% ----- gris de la fase administrativa -----
\definecolor{adminG}{HTML}{8A98A6}
% \gbar{fila}{semana inicial}{semana final}{color}{etiqueta}
\newcommand{\gbar}[5]{%
  \fill[#4,rounded corners=2.2pt] ({#2-1},{-#1+1-0.20}) rectangle ({#3},{-#1+1+0.20});%
  \node[anchor=west,align=left,text width=4.9cm,font=\scriptsize,text=black!80]
       at (-8.2,{-#1+1}) {#5};}

\begin{table}[H]
\caption{Cronograma de actividades (octubre 2026 -- septiembre 2027), por meses.}
\label{tab:cronograma}

% ----- leyenda de fases (alineada a la izquierda) -----
\noindent{\scriptsize
\begin{tikzpicture}[baseline]
  \fill[adminG,rounded corners=1.5pt](0,0) rectangle (0.32,0.26);
  \node[anchor=west] at (0.4,0.13) {Administrativo};
  \fill[g3,rounded corners=1.5pt](3.0,0) rectangle (3.32,0.26);
  \node[anchor=west] at (3.4,0.13) {Desarrollo (ASR, MT, TTS)};
  \fill[g10,rounded corners=1.5pt](7.6,0) rectangle (7.92,0.26);
  \node[anchor=west] at (8.0,0.13) {Investigaci{\'o}n y cierre};
\end{tikzpicture}}

\smallskip
\noindent\resizebox{\linewidth}{!}{%
\begin{tikzpicture}[x=0.92cm,y=0.82cm,font=\scriptsize]
  % --- banda de meses (navy): 12 meses, oct 2026 -> sep 2027 ---
  \fill[g5,rounded corners=1.5pt] (0,0.74) rectangle (12,1.6);
  \foreach \m/\pos in {Oct/0.5,Nov/1.5,Dic/2.5,Ene/3.5,Feb/4.5,Mar/5.5,Abr/6.5,May/7.5,Jun/8.5,Jul/9.5,Ago/10.5,Sep/11.5}
     \node[white,font=\bfseries\scriptsize] at (\pos,1.17) {\m};
  % --- numeros de mes (1..12) ---
  \foreach \w in {1,...,12}
     \node[font=\tiny,text=black!55] at ({\w-0.5},0.45) {\w};
  % --- separadores de mes ---
  \foreach \xg in {0,...,12}
     \draw[black!18,very thin] (\xg,0.74) -- (\xg,-11.6);
  % --- lineas horizontales entre filas ---
  \foreach \r in {0,...,12}
     \draw[black!10,very thin] (0,{-\r+0.5}) -- (12,{-\r+0.5});
  % --- barras (fila, mes inicial, mes final, color, etiqueta) ---
  \gbar{1}{1}{1}{adminG}{Aprobaci{\'o}n del plan de tesis}
  \gbar{2}{1}{3}{g1}{Revisi{\'o}n del estado del arte}
  \gbar{3}{2}{5}{g2}{Recolecci{\'o}n y preparaci{\'o}n del corpus (salud)}
  \gbar{4}{4}{6}{g3}{M{\'o}dulo ASR (voz $\to$ texto)}
  \gbar{5}{5}{7}{g4}{M{\'o}dulo MT (traducci{\'o}n)}
  \gbar{6}{6}{8}{g5}{M{\'o}dulo TTS (texto $\to$ voz)}
  \gbar{7}{8}{9}{g6}{Integraci{\'o}n del \emph{pipeline}}
  \gbar{8}{9}{10}{g8}{Evaluaci{\'o}n experimental (WER, BLEU, MOS)}
  \gbar{9}{10}{11}{g9}{Validaci{\'o}n m{\'e}dico-paciente (simulada)}
  \gbar{10}{11}{11}{g10}{An{\'a}lisis de resultados}
  \gbar{11}{10}{12}{g11}{Redacci{\'o}n del informe final}
  \gbar{12}{12}{12}{g12}{Sustentaci{\'o}n}
\end{tikzpicture}}
\notaAPA{Elaboraci{\'o}n propia. Cada columna representa un mes (octubre 2026 -- septiembre 2027); el color identifica la fase y el solapamiento de las barras refleja el desarrollo parcialmente paralelo de los m{\'o}dulos ASR, MT y TTS.}
\end{table}

\subsection{Presupuesto y financiamiento}
La Tabla~\ref{tab:presupuesto} resume los recursos estimados para el desarrollo del proyecto.

\begin{table}[H]
\centering
\caption{Recursos para el desarrollo del proyecto de tesis (soles).}
\label{tab:presupuesto}
\small
\begin{tabularx}{\linewidth}{@{}l X r@{}}
\toprule\hdr{Tipo} & \hdr{Descripci{\'o}n} & \hdr{Costo (S/)} \\
\midrule
Hardware & Laptop/PC con GPU (entrenamiento y \emph{fine-tuning}) & 4\,500.00 \\
C{\'o}mputo & Cr{\'e}ditos de GPU en la nube (horas de ajuste) & 900.00 \\
Datos & Grabaci{\'o}n y validaci{\'o}n del corpus cl{\'i}nico (incentivos a hablantes) & 800.00 \\
Administrativos & Impresiones, inscripci{\'o}n y tr{\'a}mites de tesis & 600.00 \\
Movilidad & Traslados al centro de salud y a comunidades & 300.00 \\
Otros & Gastos no previstos & 250.00 \\
\multicolumn{2}{@{}l}{\textbf{Total}} & \textbf{7\,350.00} \\
\bottomrule
\end{tabularx}
\notaAPA{Montos estimados en soles (S/). Elaboraci{\'o}n propia.}
\end{table}

\subsection{Equipo de investigaci{\'o}n}
La Tabla~\ref{tab:equipo} detalla el equipo responsable del proyecto.

\begin{table}[H]
\centering
\caption{Equipo de investigaci{\'o}n.}
\label{tab:equipo}
\small
\begin{tabularx}{\linewidth}{@{}l X@{}}
\toprule\hdr{Rol} & \hdr{Responsable / funci{\'o}n} \\
\midrule
Tesista   & Alain Armando Condori Contreras ---dise{\~n}o, desarrollo y evaluaci{\'o}n del sistema---. \\
Asesor    & Mtro. Ing. Ra{\'u}l Huillca Huallparimachi ---orientaci{\'o}n metodol{\'o}gica y t{\'e}cnica---. \\
Colaboradores & Personal de salud biling{\"u}e (validaci{\'o}n del corpus) y participantes voluntarios (consultas simuladas). \\
\bottomrule
\end{tabularx}
\notaAPA{Elaboraci{\'o}n propia. Los colaboradores intervienen en la validaci{\'o}n del corpus y en la actuaci{\'o}n de las consultas simuladas.}
\end{table}
    % Cap. V  - Aspectos administrativos

% ---- Cierre ----
% ================= BIBLIOGRAFÍA =================
\newpage
\defbibnote{vigencia}{\footnotesize\textit{Nota sobre la vigencia de las fuentes.}
Se priorizaron publicaciones de los {\'u}ltimos cinco a{\~n}os (2021--2026). Las referencias
anteriores se conservan por ser (a)~obras \emph{seminales} que definen los m{\'e}todos empleados
---el \emph{Transformer} \parencite{vaswani2017attention}, BLEU \parencite{papineni2002bleu} y chrF
\parencite{popovic2015chrf}---; (b)~\emph{libros} de metodolog{\'i}a en su edici{\'o}n vigente
\parencite{hernandez2018metodologia}; y (c)~\emph{fuentes primarias oficiales} o trabajos
fundacionales del PLN para el quechua ---los Censos Nacionales \parencite{inei2017} y los primeros
sistemas de traducci{\'o}n \parencite{ortega2020quechua} y reconocimiento de voz
\parencite{Zevallos202061} para esta lengua---.\par\vspace{4pt}}
\printbibliography[heading=bibintoc,title={Referencias},prenote=vigencia]
    % Referencias (APA7)
% ================= ANEXOS =================
\newpage
\section*{Anexos}
\addcontentsline{toc}{section}{Anexos}

% ----- Anexo 1 -----
\subsection*{Anexo 1: Ficha t{\'e}cnica de registro experimental}
\addcontentsline{toc}{subsection}{Anexo 1: Ficha t{\'e}cnica de registro experimental}
Plantilla para registrar cada \emph{turno de traducci{\'o}n} durante las consultas simuladas.

\begin{table}[H]
\centering
\caption{Ficha t{\'e}cnica de registro por turno (plantilla).}
\label{tab:ficha}
\small
\begin{tabularx}{\linewidth}{@{}l X@{}}
\toprule\hdr{Campo} & \hdr{Valor registrado} \\
\midrule
ID de turno / sesi{\'o}n            & \rule{5cm}{0.3pt} \\
Escenario (triaje/s{\'i}ntomas/indicaciones) & \rule{5cm}{0.3pt} \\
Direcci{\'o}n (\sig{quz}$\to$es / es$\to$\sig{quz}) & \rule{5cm}{0.3pt} \\
Transcripci{\'o}n ASR / WER         & \rule{5cm}{0.3pt} \\
Traducci{\'o}n MT / BLEU / chrF     & \rule{5cm}{0.3pt} \\
Calidad de voz (MOS)               & \rule{5cm}{0.3pt} \\
Latencia por turno (s)             & \rule{5cm}{0.3pt} \\
Observaciones                      & \rule{5cm}{0.3pt} \\
\bottomrule
\end{tabularx}
\notaAPA{Elaboraci{\'o}n propia. Se completa una ficha por cada turno de traducci{\'o}n registrado.}
\end{table}

% ----- Anexo 2 -----
\subsection*{Anexo 2: Plantilla del banco de frases de consulta (evaluaci{\'o}n)}
\addcontentsline{toc}{subsection}{Anexo 2: Plantilla del banco de frases de consulta (evaluaci{\'o}n)}
Gui{\'o}n base de frases por escenario; el equivalente en quechua de Cusco se valida con hablantes
nativos antes de su uso.

\begin{table}[H]
\centering
\caption{Banco de frases de consulta simulada (ejemplos; \sig{quz} por validar).}
\label{tab:banco}
\small
\begin{tabularx}{\linewidth}{@{}l X l@{}}
\toprule\hdr{Categor{\'i}a} & \hdr{Enunciado (espa{\~n}ol)} & \hdr{Equivalente (\sig{quz})} \\
\midrule
Triaje       & {\textquestiondown}Qu{\'e} lo trae a la consulta?          & [por validar] \\
Triaje       & {\textquestiondown}Hace cu{\'a}nto tiempo se siente mal?    & [por validar] \\
S{\'i}ntomas & {\textquestiondown}D{\'o}nde le duele?                      & [por validar] \\
S{\'i}ntomas & {\textquestiondown}Tiene fiebre o escalofr{\'i}os?          & [por validar] \\
Indicaciones & Tome esta pastilla cada 8 horas.                          & [por validar] \\
Indicaciones & Regrese en 3 d{\'i}as si no mejora.                        & [por validar] \\
\bottomrule
\end{tabularx}
\notaAPA{Elaboraci{\'o}n propia. El equivalente en quechua de Cusco (\sig{quz}) se valida con hablantes nativos antes de su uso.}
\end{table}

% ----- Anexo 3 -----
\subsection*{Anexo 3: Cuestionario de percepci{\'o}n (escala Likert)}
\addcontentsline{toc}{subsection}{Anexo 3: Cuestionario de percepci{\'o}n (escala Likert)}
Escala de 1 (\emph{totalmente en desacuerdo}) a 5 (\emph{totalmente de acuerdo}); se aplica al
personal y a los pacientes-actores.

\begin{table}[H]
\centering
\caption{{\'I}tems del cuestionario de percepci{\'o}n.}
\label{tab:likert}
\small
\begin{tabularx}{\linewidth}{@{}c X c@{}}
\toprule\hdr{N.\textordmasculine} & \hdr{{\'I}tem} & \hdr{1--5} \\
\midrule
1 & El sistema facilit{\'o} la comunicaci{\'o}n durante la consulta.     & \rule{1.6cm}{0.3pt} \\
2 & Conf{\'i}o en la traducci{\'o}n que entreg{\'o} el sistema.          & \rule{1.6cm}{0.3pt} \\
3 & Comprend{\'i} con claridad las indicaciones m{\'e}dicas.            & \rule{1.6cm}{0.3pt} \\
4 & La voz sintetizada fue clara y natural.                             & \rule{1.6cm}{0.3pt} \\
5 & El tiempo de respuesta del sistema fue adecuado.                    & \rule{1.6cm}{0.3pt} \\
6 & Usar{\'i}a el sistema en una consulta real.                         & \rule{1.6cm}{0.3pt} \\
\bottomrule
\end{tabularx}
\notaAPA{Elaboraci{\'o}n propia. Escala Likert de 1 (totalmente en desacuerdo) a 5 (totalmente de acuerdo).}
\end{table}

% ----- Anexo 4 -----
\subsection*{Anexo 4: Modelo de consentimiento informado}
\addcontentsline{toc}{subsection}{Anexo 4: Modelo de consentimiento informado}
\noindent\textit{Yo, \rule{5cm}{0.3pt}, declaro haber sido informado(a) sobre los objetivos del
estudio ``\nuevo{Desarrollo de un sistema bidireccional de traducci{\'o}n de voz Quechua--Espa{\~n}ol
mediante aprendizaje profundo para consultas m{\'e}dicas en Cusco}''. Entiendo que participar{\'e} de forma voluntaria en consultas
\textbf{simuladas} ---sin datos cl{\'i}nicos reales---, que las grabaciones ser{\'a}n anonimizadas y
usadas solo con fines de investigaci{\'o}n, y que puedo retirarme en cualquier momento sin
perjuicio.}

\vspace{1.2cm}
\noindent\rule{6cm}{0.3pt}\hfill\rule{4cm}{0.3pt}\\
\noindent Firma del participante \hfill Fecha

          % Anexos

\end{document}
